\documentclass[11pt,a4paper]{ctexart}
\usepackage[margin=2.25cm]{geometry}
\usepackage{amsmath,amssymb,mathtools}
\usepackage{booktabs,longtable,tabularx,array}
\usepackage{enumitem}
\usepackage{xcolor}
\usepackage[hidelinks]{hyperref}
\usepackage{microtype}

\definecolor{deepblue}{RGB}{32,70,118}
\definecolor{softgray}{RGB}{245,246,248}
\newcommand{\R}{\mathbb R}
\newcommand{\C}{\mathbb C}
\newcommand{\N}{\mathbb N}
\newcommand{\eps}{\varepsilon}
\newcommand{\dd}{\,\mathrm d}
\newcommand{\ess}{\operatorname*{ess}}
\newcommand{\Var}{\operatorname{Var}}
\newcommand{\supp}{\operatorname{supp}}
\newcommand{\dist}{\operatorname{dist}}
\newcommand{\Lip}{\operatorname{Lip}}
\newcommand{\osc}{\operatorname{osc}}
\newcommand{\diam}{\operatorname{diam}}
\newcommand{\id}{\mathrm I}
\newcommand{\weakto}{\rightharpoonup}
\newcommand{\planitem}[1]{\par\smallskip\noindent\textcolor{deepblue}{\textbf{#1}}\quad}
\newcommand{\unithead}[2]{\subsection{#1\quad #2}}

\setlist[itemize]{leftmargin=2em,itemsep=0.25em,topsep=0.35em}
\setlist[enumerate]{leftmargin=2.2em,itemsep=0.3em,topsep=0.35em}
\renewcommand{\arraystretch}{1.2}

\title{《有限观测实分析》第四章修订计划（v2）\\
\large 局部平均、累计变化、弱测试与时间演化}
\author{}
\date{}

\begin{document}
\maketitle
\tableofcontents

\section{修订目标与实施结论}

第四章围绕同一个识别问题展开：\emph{当对象只能通过局部平均、区间增量、测试函数或有限时间轨道被观察时，怎样以尺度一致的估计稳定这些观测，并从极限中识别点值、密度、弱导数和生成元？}
全章依次推进四条链：
\[
\begin{aligned}
&\text{尺度链：局部平均}\longrightarrow\text{覆盖选择}\longrightarrow
  \text{极大控制}\longrightarrow\text{函数与测度的微分},\\
&\text{累计链：区间增量}\longrightarrow\text{总变差}\longrightarrow
  \text{Stieltjes 测度}\longrightarrow\text{普通导数所能识别的部分},\\
&\text{测试链：平移与卷积}\longrightarrow\text{光滑化}\longrightarrow
  \text{分布导数}\longrightarrow\text{Sobolev 空间与紧嵌入},\\
&\text{演化链：核的复合}\longrightarrow\text{近似恒等}\longrightarrow
  C_0\text{-半群}\longrightarrow\text{生成元与预解式}.
\end{aligned}
\]
全章以 $m$ 表示 Lebesgue 测度、$m^*$ 表示其外测度，$\omega_n=|B(0,1)|$ 表示单位球体积，
$\mathbf1_E$ 表示集合 $E$ 的示性函数，$\delta_x$ 表示点 $x$ 的 Dirac 测度；
$C_b$ 表示有界连续函数，$C_0$ 表示在无穷远处消失的连续函数，
$A\asymp B$ 表示两者之比由与尺度无关的正数上下控制。
记 $K\Subset\Omega$ 表示 $K$ 为紧集且 $K\subset\Omega$；$\ess\supp f$ 是补集上
$f=0$ a.e. 的最大开集的补集，$\supp f$ 表示普通拓扑支撑。
函数 $u$ 称为 Lipschitz，若存在 $L<\infty$ 使
$|u(x)-u(y)|\le L|x-y|$；符号 $I$ 按上下文表示单位矩阵或恒等算子；标准 Gaussian 测度 $\gamma_I$ 由密度
$(2\pi)^{-n/2}e^{-|x|^2/2}$ 定义；对半正定矩阵 $Q$ 取其对称平方根
$Q^{1/2}$，并记 $\gamma_Q=(Q^{1/2})_*\gamma_I$（4.4.1 证明这一定义与任意因子
$Q=AA^{\mathsf T}$ 给出的推前一致）。
空间函数的梯度与 Laplacian 统一记为
$\nabla u=(\partial_1u,\ldots,\partial_nu)$、
$\Delta u=\sum_{j=1}^n\partial_j^2u$（$\partial_j=\partial/\partial x_j$）；当 $F$ 是区间上的 BV 函数时，
$F(x-)$ 表示左极限，并记内点跳幅
$\Delta F(x)=F(x)-F(x-)$，与空间函数的 $\Delta$ 由上下文区分。
Heaviside 函数记为 $H=\mathbf1_{(0,\infty)}$（取 $H(0)=0$），并约定
$\operatorname{sgn}x=1$（$x>0$）、$-1$（$x<0$）、$0$（$x=0$）。
记 $\ell^2(\mathbb Z)=\{a=(a_k)_{k\in\mathbb Z}:\sum_{k\in\mathbb Z}|a_k|^2<\infty\}$，
范数为 $\|a\|_2=(\sum_k|a_k|^2)^{1/2}$。
对赋范空间写 $X\hookrightarrow Y$ 表示连续嵌入，$X\Subset Y$ 表示紧嵌入
（$X$ 的有界集在 $Y$ 中相对紧）。
对有限有符号或复测度 $\mu$，$|\mu|$ 为全变差测度，
$\|\mu\|_{TV}=|\mu|(\mathbb R^n)$（在区间上相应限制）；$C([a,b])^*$ 表示
$C([a,b])$ 的连续对偶。
球面 $\mathbb S^{n-1}$ 取标准表面积测度，极坐标公式中的径向积分均按第三章的约定理解。
后文出现的 Beta 函数预先记为
\[
B(\alpha,\beta)=\int_0^1s^{\alpha-1}(1-s)^{\beta-1}\,\dd s
 =\frac{\Gamma(\alpha)\Gamma(\beta)}{\Gamma(\alpha+\beta)}
 \qquad(\alpha,\beta>0),
\]
其中 $\Gamma(\alpha)=\int_0^\infty t^{\alpha-1}e^{-t}\,\dd t$ 为 Euler Gamma 函数。
平移与卷积统一约定为
$\tau_hf(x)=f(x-h)$、
$(f*g)(x)=\int_{\R^n}f(y)g(x-y)\,\dd y$（在相应可积性条件下）。
全空间中心极大函数预先记为
$Mf(x)=\sup_{r>0}|B(x,r)|^{-1}\int_{B(x,r)}|f|\,\dd m$；局部版本在 4.1.2 中另行限定尺度。
记 $e_1$ 为第一坐标向量。
所有 $r\downarrow0$ 的量词均按有向集
$D=(0,r_0)$、$r\preceq s\Longleftrightarrow r\ge s$ 理解；正有理半径在允许范围内共尾。
四条链共用同一证明节奏：先在连续、光滑或有限维对象上完成计算；再建立与尺度无关的估计；随后用前三章的稠密性、正则性、完备性或受控收敛传递结论；最后用点值、测度、测试函数或生成元识别极限。

全章设置十九个正文单元，章幅建议为 216--223 页（按下表区间相加为 216--223 页）。扩展篇幅主要放在覆盖选择的量化步骤、相对测度微分、BV 端点规范与 indicatrix、Sobolev 几何假设、紧性反例、核解方程以及半群预解式的双向证明。每个正式编号的结论都在此前获得定义和依赖；每个例题都完成一个可计算或可证明的问题，跨单元例题标明回访位置，诊断例则逐项检验定理假设。

\subsection{篇幅与节奏总表}

\begin{longtable}{p{0.14\textwidth}p{0.56\textwidth}p{0.14\textwidth}}
\toprule
单元 & 核心交付 & 建议页数\\
\midrule
章首 & 四条链、前三章接口、例题/诊断账本与符号约定 & 6\\
4.1.1 & 可允许的局部平均、振荡可测性、覆盖选择 & 8\\
4.1.2 & 极大弱型、弱 $L^1$、Marcinkiewicz 特例 & 10\\
4.1.3 & Lebesgue 点、密度点、精确代表 & 10\\
4.1.4 & 有限 Radon 测度的微分与相对密度 & 14\\
4.2.1 & BV、Jordan 分解与端点完整的测度编码 & 11\\
4.2.2 & 单调函数可微与 Stieltjes 测试 & 11--12\\
4.2.3 & AC、FTC、换元、曲线和复合边界 & 11\\
4.2.4 & BV 三分解、indicatrix、面积公式与 Banach--Zarecki & 15--16\\
4.3.1 & Young 不等式、mollifier 与光滑稠密 & 11\\
4.3.2 & 分布、分布导数与弱极限 & 10--11\\
4.3.3 & Sobolev 基础、一维代表、ACL、Meyers--Serrin & 12\\
4.3.4 & Poincar\'e、Friedrichs 与 Lipschitz 延拓 & 12\\
4.3.5 & Sobolev--Morrey 嵌入及端点诊断 & 11--12\\
4.3.6 & Kolmogorov--Riesz 与 Rellich 紧嵌入 & 10--11\\
4.4.1 & 卷积代数、测度卷积与有限测度平滑 & 10\\
4.4.2 & 近似恒等的范数/点态恢复与失败诊断 & 9--10\\
4.4.3 & 热核、Poisson 核及正时间方程 & 12--13\\
4.4.4 & 矩阵实验、$C_0$-半群与生成元 & 11\\
4.4.5 & 轨道微分、能量、Laplace--Bochner 预解式 & 12\\
\bottomrule
\end{longtable}

\subsection{前三章的输入与第四章的新工作}

正文首次引用前以“输入卡”列出下表中的正式结果，并统一使用前三章源码中的实际定理标签或定理名称。

\begin{longtable}{p{0.13\textwidth}p{0.40\textwidth}p{0.38\textwidth}}
\toprule
来源 & 已建立的输入 & 在第四章中的落点\\
\midrule
第一章 & 有向性与共尾、紧性及有限子覆盖、Euclidean 第二可数性、完备性、$C_0$ 的完备性与 $C_c$ 在 $C_0$ 中稠密、距离函数分离、光滑 bump、等度连续与 Arzel\`a--Ascoli、Baire 定理及其标量一致有界原型 & 覆盖族的有限化；Stieltjes 泛函范数的连续分离函数；Rellich 的有限网；本章用 Baire 推出算子族与半群的局部一致有界\\
第二章 & Lebesgue/Radon 正则性、Hahn--Jordan 与全变差、Lebesgue 分解、Radon--Nikodym、乘积与推前测度、Lebesgue--Stieltjes 增量及左端原子、一致捕获质量的紧集条件 & 测度微分；BV--测度对应；加法推前定义测度卷积；非紧空间上的质量控制\\
第三章 & $L^1_{\rm loc}$、Fatou/DCT/MCT/Vitali、层蛋糕与 Chebyshev、Tonelli--Fubini、H\"older 与 Minkowski、$L^p$ 完备性与对偶、$C_c$ 稠密、$L^p$ 平移连续（标签 \texttt{prop:3-3-3-translation-continuity}）、$L^1$ 卷积的可测性和估计、AC 的积分表示、Bochner 积分 & 极大函数的分布估计；稠密推广；卷积光滑化；一维 Sobolev 代表；半群轨道积分与预解式\\
\bottomrule
\end{longtable}

第四章的新内容由这些接口自然产生：覆盖引理把全体尺度压缩成可加估计；极大函数把尺度一致控制转成几乎处处恢复；测度微分识别 Stieltjes 增量的绝对连续部分；测试函数把经典导数扩张为分布导数；平移模与尾部控制刻画函数族的强紧性；核的时间复合最终抽象为 $C_0$-半群和生成元。

\subsection{贯穿例题与诊断例账本}

以下对象沿四条主线承担明确的计算、反例诊断或跨单元回访。

\begin{longtable}{p{0.09\textwidth}p{0.43\textwidth}p{0.40\textwidth}}
\toprule
编号 & 对象与具体任务 & 承担的问题与回访\\
\midrule
E1 & $\mathbf1_{B(0,1)}$ 的中心极大函数；分区计算远处衰减 $M\mathbf1_{B(0,1)}(x)\asymp |x|^{-n}$ & 4.1.2 检验弱 $(1,1)$ 的端点；4.4.2 回访极大控制下的核恢复\\
E2 & 在 $\R$ 上取交替环层集合，精确计算两列区间平均为 $2/3$ 与 $1/3$ & 4.1.1 展示全尺度问题；4.1.3 判定密度失败；4.4.2 比较偏心核\\
E3 & 从 $[0,1]$ 开始构造 Smith--Volterra 集 $K_{SV}$：第 $m$ 步从每个剩余区间删除长度 $4^{-m}$ 的中央开区间，计算 $m(K_{SV})=1/2$、$\operatorname{int}K_{SV}=\varnothing$、$\partial K_{SV}=K_{SV}$ & 4.1.3 对照拓扑边界与几乎处处密度，并把每个密度窗口落实为可计算比例\\
E4 & 在 $[0,1]$ 上取 $p\delta_{1/3}+q\mu_C+(1-p-q)m|_{[0,1]}$，其中 $\mu_C$ 是第二章由二进制柱集推前（等价地由 Carath\'eodory 扩张）得到、每个第 $m$ 级基本区间质量为 $2^{-m}$ 的标准三分 Cantor 概率测度、$p,q>0$、$p+q<1$ & 测度密度、BV 的跳跃/奇异连续/绝对连续三分解\\
E5 & 互异跳点的有限阶梯函数 & 变差、Stieltjes 测度、积分范数、FTC 失效位置\\
E6 & $[0,1]$ 上的 $x^2\sin(x^{-2})$（在零点补零）；沿显式交错取值列求变差下界 & 4.2.1 诊断连续、a.e. 可微与 BV 之间的严格边界\\
E7 & $[0,1]$ 上的 Cantor 函数 & 导数遗漏的奇异变化、$N$ 性质与 Banach--Zarecki\\
E8 & 在原点置零的 $|x|^{-1/2}\mathbf1_{(-1,1)}$ 的自卷积；在原点直接算出发散积分 & 4.3.1 检验 $L^1*L^1$ 的“几乎处处有限”边界\\
E9 & 区间示性函数的三角形卷积与标准 mollifier & 支撑、连续性、光滑化、边界过渡层\\
E10 & Heaviside 函数和 $|x|/2$ & $H'=\delta_0$、$(|x|/2)''=\delta_0$、一维弱导数\\
E11 & 归一化尖峰 $k^{n/p}\mathbf1_{B(0,1/k)}$ 与平移列 $\eta(\cdot-ke_1)$（取 $0\ne\eta\in C_c^\infty(B(0,1))$）；逐项核对紧性判据 & 4.3.6 分离平移模与尾部条件的职责\\
E12 & 高斯核、Poisson 核与 $\ell^2(\mathbb Z)$ 上的对角算子族 $a\mapsto(e^{-tk^2}a_k)_k$ & 4.4.3--4.4.5 依次检验近似恒等、半群律、无界生成元与预解式\\
\bottomrule
\end{longtable}

\subsection{全章符号与首次定义纪律}

\begin{enumerate}
\item 全章以 $m$ 表示 Lebesgue 测度，$\omega_n=|B(0,1)|$ 表示单位球体积，
      $\mathbf1_E$ 表示集合 $E$ 的示性函数；对连续 $g$ 在紧区间上的连续模记为
      $\omega_g(\delta)=\sup\{|g(x)-g(y)|:|x-y|\le\delta\}$。
      $C_b$ 表示有界连续函数，$C_0$ 表示在无穷远处消失的连续函数。
\item $L^p$ 元素需要点值时先把异常零测集上的值置零，固定处处有限的可测代表；精确代表在 Lebesgue 微分定理之后定义。
\item $f\in L^1_{\rm loc}(\Omega)$ 的球平均只取 $0<r<\dist(x,\Omega^c)$，并约定
$\dist(x,\varnothing)=\infty$；进入全空间卷积时再使用全空间函数或明确的延拓。
\item 弱 $L^p$ 准范数、分布函数、弱收敛的测试类、BV 范数、左端原子、测试函数空间的收敛、Sobolev 域的几何条件、半群生成元定义域，均在第一次出现在公式前定义。
\item 对 $L^p$ 等价类使用 $\ess\supp f$；对普通函数使用拓扑支撑。写 $u_\Omega$ 时同时声明 $0<|\Omega|<\infty$ 与 $u\in L^1(\Omega)$。
\item 复 Banach/Hilbert 空间中的能量公式取内积实部；半群增长界统一写成 $\|T_t\|\le Me^{\omega t}$。
\end{enumerate}

\section{逐单元写作计划}

\subsection*{4.1　尺度链：从全部小球到几乎处处恢复}

\unithead{4.1.1}{局部平均、振荡与覆盖选择}

\planitem{驱动问题}
给定 $f\in L^1_{\rm loc}(\Omega)$，沿一列半径观察 $A_{r_k}f(x)$ 为什么不足以确定 $r\downarrow0$ 的极限？若每个坏点都选一个见证坏行为的球，怎样从大量重叠球中抽出可求和的信息？

\planitem{定义与正式结论}
对 $x\in\Omega$ 和 $0<r<\dist(x,\Omega^c)$ 定义
\[
A_rf(x)=\frac1{\omega_n r^n}\int_{B(x,r)}f(y)\dd y.
\]
固定可测代表后定义
\[
 \Omega_f(x)=\limsup_{\substack{r\downarrow0,\ r\in\mathbb Q\\
0<r<\dist(x,\Omega^c)}}
 \frac1{|B(x,r)|}\int_{B(x,r)}|f(y)-f(x)|\dd y.
\]
称球族在集合 $E$ 上细覆盖（Vitali 覆盖），是指每个 $x\in E$ 都是族中球的中心，
且对任意 $\delta>0$ 都有以 $x$ 为中心、半径小于 $\delta$ 的球；4.1.4 的细球族是
这一概念的有界版本。
正文依次证明：
\begin{enumerate}
\item 固定允许半径时，$x\mapsto\int_{B(x,r)}|f|$ 在内缩开集上连续；固定 $x$ 时，归一化半径积分连续。因此正有理半径在允许半径集内共尾，公式中的可数上极限等于沿全部实半径的上极限。
\item $(x,y)\mapsto\mathbf1_{\{|x-y|<r\}}|f(y)-f(x)|$ 联合可测；由参数积分可测性及有理半径的可数性，$\Omega_f$ 可测。
\item 若 $Q(x,r)=x+[-r,r]^n$，则
$B(x,r)\subset Q(x,r)\subset B(x,\sqrt n r)$。对非负 $f$（对一般 $f$ 应用于 $|f|$），若
$A_r^\square f=(2r)^{-n}\int_{Q(x,r)}f$，则写出双向常数
\[
A_rf(x)\le\frac{2^n}{\omega_n}A_r^\square f(x),\qquad
A_r^\square f(x)\le\frac{\omega_n n^{n/2}}{2^n}A_{\sqrt n r}f(x).
\]
\item $5r$ 覆盖引理：半径有上界的欧氏球族可抽取至多可数个两两不交球 $B_j$，使原球族的并包含于 $\bigcup_j5B_j$。证明按二进尺度分层，每层作极大不交选择，并以第二可数性给出可数性。
\end{enumerate}

\planitem{证明节奏}
先计算平移球的对称差并用积分绝对连续性处理中心变化；再将振荡写成二变量积分完成可测性；随后用两个相交区间展示“直接取不交子族会漏点”，由二进尺度选择解释放大因子 $5$ 的来源。覆盖证明中的每个有限阶段都写出已选择球、待选球及半径比较，最后才取可数并。

\planitem{主例}
\begin{itemize}
\item 在 $\R$ 上令 $f(0)=0$、$f(x)=|x|^{-1/2}$（$x\ne0$），则 $A_rf(0)=2r^{-1/2}\to+\infty$，说明局部可积只保证积分有限，不保证平均有有限极限。
\item 仍在 $\R$ 上，令
\[
E=\bigcup_{j\ge0}\{x:2^{-2j-1}<|x|\le2^{-2j}\}.
\]
完整算出 $A_{2^{-2m}}\mathbf1_E(0)=2/3$ 与 $A_{2^{-2m-1}}\mathbf1_E(0)=1/3$。
\item 取 $I_j=(-j\eps,1+(M-j)\eps)$，$j=1,\ldots,M$，并令 $f=\mathbf1_{(0,1)}$。除计算 $(0,1)$ 上覆盖重数为 $M$ 外，进一步核对
\[
\sum_{j=1}^M\int_{I_j}f=M,
\qquad
\int_{\bigcup_jI_j}f=1.
\]
随后对这族区间实际执行不交选择并比较膨胀覆盖，使去重步骤直接对应弱型证明中的可加估计。
\end{itemize}

\planitem{输出与练习}
本单元输出允许半径、可测振荡和 $5r$ 选择三项工具。练习分别要求把球改成立方体、核对有限族可取 $3r$ 的精确条件、构造另一组具有两个子列密度的径向集合。

\unithead{4.1.2}{Hardy--Littlewood 极大控制与插值}

\planitem{驱动问题}
坏点各自选出的半径不同，不能由一个固定半径的积分估计统一处理。定义全尺度上确界后，要证明的问题是
\[
m\{x:Mf(x)>\lambda\}\quad\text{能否只由}\quad \|f\|_1/\lambda
\quad\text{控制？}
\]

\planitem{定义与正式结论}
在 $\R^n$ 上定义中心极大函数
\[
Mf(x)=\sup_{r>0}\frac1{|B(x,r)|}\int_{B(x,r)}|f(y)|\dd y,
\]
并用有理半径和上一单元的连续性证明可测。对可测 $g$ 定义分布函数
$d_g(\lambda)=m\{|g|>\lambda\}$；当 $1\le p<\infty$ 时再定义
$\|g\|_{L^{p,\infty}}=\sup_{\lambda>0}\lambda d_g(\lambda)^{1/p}$。
称算子 $T$ 次线性，是指 $|T(f+g)|\le |Tf|+|Tg|$ 且
$|T(\alpha f)|=|\alpha|\,|Tf|$（逐点或几乎处处）。称 $T$ 为弱 $(p,p)$ 型，若存在
$C$ 使 $\|Tf\|_{L^{p,\infty}}\le C\|f\|_p$（$p=1$ 时按分布函数的等价写法理解）；称为强
$(p,p)$ 型，若 $\|Tf\|_p\le C\|f\|_p$。下文的弱型与强型估计均对其可测输出理解。
正式建立：
\begin{enumerate}
\item Hardy--Littlewood 弱 $(1,1)$ 估计
\[
m\{Mf>\lambda\}\le \frac{5^n}{\lambda}\|f\|_1,
\qquad f\in L^1(\R^n),\ \lambda>0.
\]
\item $M$ 的 $L^\infty$ 上界 $Mf\le\|f\|_\infty$；等号的范数结论留到微分定理之后作为推论。
\item Marcinkiewicz 特例：次线性算子若为弱 $(1,1)$ 且强 $(\infty,\infty)$，则对 $1<p<\infty$ 为强 $(p,p)$。证明用阈值分解 $f=f_{\le a\lambda}+f_{>a\lambda}$、层蛋糕公式和 Tonelli，完整计算 $p/(p-1)$ 型常数。
\item 对 $R>0$ 定义局部极大函数
\[
M_Rf(x)=\sup_{0<r<R}\frac1{|B(x,r)|}\int_{B(x,r)}|f(y)|\dd y.
\]
证明 $M_R$ 的相同弱型估计以及局部化到相对紧开集的版本。
\end{enumerate}

\planitem{证明节奏}
对每个 $x\in\{Mf>\lambda\}$ 选取见证球；先在紧子集上用有限覆盖，再以 $5r$ 引理抽取不交球，写出
\[
\lambda\sum_j|B_j|<\sum_j\int_{B_j}|f|\le\|f\|_1,
\qquad
m\Bigl(\bigcup_j5B_j\Bigr)\le5^n\sum_j|B_j|.
\]
随后用开集的内逼近撤去紧子集。插值证明先用一页固定分布函数、层蛋糕与截断记号，随后连续完成截断积分。

\planitem{主例}
\begin{itemize}
\item 对 $f=\mathbf1_{B(0,1)}$，分别处理 $|x|\le2$ 与 $|x|>2$，得到远处 $Mf(x)\asymp |x|^{-n}$；由此看见 $Mf\notin L^1$ 而 $Mf\in L^{1,\infty}$。
\item 在 $(0,1)$ 上取 $g_1(x)=1/x$、$g_2(x)=1/(1-x)$。计算两者弱 $L^1$ 准范数均为 $1$，而 $g_1+g_2=1/[x(1-x)]$ 的弱准范数为 $4$，明确弱 $L^1$ 的准三角边界。
\item 枚举 $\mathbb Q^n=\{q_k\}$，取
\[
f(x)=\sum_{k\ge1}2^{-k}|x-q_k|^{-n/2}
\mathbf1_{B(q_k,1)}(x).
\]
先把所得和在可数集 $\mathbb Q^n$ 上的值置零；再用极坐标逐项算出 $\|f\|_1<\infty$，并由第 $k$ 项在 $q_k$ 的球平均按 $r^{-n/2}$ 发散推出 $Mf(q_k)=+\infty$。这说明弱型结论允许无穷值集稠密，同时强制其为零测。
\end{itemize}

\planitem{输出与练习}
输出弱型、强 $L^p$ 和局部极大估计。练习要求核算非中心极大函数与中心版本的常数比较、完成一维 rising-sun 证明的左端分量不等式、验证弱 $L^p$ 的尺度不变性。

\unithead{4.1.3}{Lebesgue 点、密度点与精确代表}

\planitem{驱动问题}
对连续函数，球平均由一致连续性恢复点值；对一般 $L^1_{\rm loc}$ 函数，怎样把连续近似的误差同时压住所有足够小的半径？

\planitem{定义与正式结论}
称 $x$ 为固定代表 $f$ 的 Lebesgue 点，若
\[
\lim_{r\downarrow0}\frac1{|B(x,r)|}
\int_{B(x,r)}|f(y)-f(x)|\dd y=0.
\]
正文完整证明：
\begin{enumerate}
\item Lebesgue 微分定理：$f\in L^1_{\rm loc}(\Omega)$ 的几乎每一点是 Lebesgue 点，且 $A_rf(x)\to f(x)$。
\item 精确代表
\[
f^*(x)=\lim_{r\downarrow0}A_rf(x)
\]
在极限存在且有限处取该极限，其余点置为 $0$（或预先固定的可测值），从而得到全域可测代表并与原函数几乎处处相等；零测修改只改变相应点的判定，不改变几乎处处结论。
\item 密度定理：对 Lebesgue 可测集 $E$，几乎每个 $x\in E$ 的密度为 $1$，几乎每个 $x\notin E$ 的密度为 $0$。
\item 有界离心率推论：若 $x\in V_k\subset B(x,r_k)$、$r_k\downarrow0$ 且 $|V_k|\ge c|B(x,r_k)|$，则在 Lebesgue 点
\[
\frac1{|V_k|}\int_{V_k}|f(y)-f(x)|\dd y\longrightarrow0.
\]
一维左右区间作为直接特例，为 4.2 的差商作准备。
\end{enumerate}

\planitem{证明节奏}
先对 $g\in C_c$ 由一致连续性证明结论；给 $f$ 选 $g$ 使 $\|f-g\|_1$ 小；把坏点包含在
\[
\{M(f-g)>\eta\}\cup\{|f-g|>\eta\},
\]
分别用弱型和 Chebyshev 控制；最后让逼近误差趋零。局部版本先用截断把待考察紧集放入一个全空间 $L^1$ 函数中。密度定理代入 $f=\mathbf1_E$ 后，立刻回访 E3 的 Smith--Volterra 集。

\planitem{主例}
\begin{itemize}
\item 对 E3 的 Smith--Volterra 集逐步求删除长度总和，得到 $m(K_{SV})=1/2$；由每个剩余区间后续仍被穿孔证明内点为空及 $\partial K_{SV}=K_{SV}$，再由密度定理得几乎每个 $K_{SV}$ 中的点具有密度 $1$。
\item 当 $n\ge2$，取球面上表面积比例为 $\theta$ 的可测区域 $\Sigma_\theta$，令 $E_\theta=\{r\omega:r>0,\omega\in\Sigma_\theta\}$；用极坐标公式算出原点在每个半径的密度均为 $\theta$。
\item 回访 E2：两列密度不同，故原点没有集合密度。
\end{itemize}

\planitem{与第三章的回扣}
第三章的 Riemann 可积判据（标签 \texttt{theorem:3-1-6-lebesgue-criterion}）在此作为一页
“局部平均回读”：推出有界 Riemann 可积函数在几乎处处由球平均恢复；本页集中以
Dirichlet 函数 $D=\mathbf1_{\mathbb Q}$ 和 Thomae 函数
$T(p/q)=1/q$（$p/q$ 为既约有理数）、$T(x)=0$（$x$ 无理）比较点态不连续与零测坏集。

\planitem{输出与练习}
输出精确代表、密度定理和偏心集合推论。练习要求证明 $\|Mf\|_\infty=\|f\|_\infty$、计算区间端点的单侧密度、把密度定理用于证明一个正测集的自差集含零点邻域。

\unithead{4.1.4}{Radon 测度的微分：从分解到局部密度}

\planitem{驱动问题}
若只知道有限有符号测度 $\nu$ 在每个小球上的质量，能否从
$\nu(B(x,r))/|B(x,r)|$ 读回 Radon--Nikodym 密度？Dirac 与 Cantor 测度的质量会在什么点态意义下消失？

\planitem{定义与正式结论}
对局部有限测度先定义球密度比
\[
D_r\nu(x)=\frac{\nu(B(x,r))}{|B(x,r)|},\qquad
\overline D\nu(x)=\limsup_{r\downarrow0}D_r\nu(x),\quad
\underline D\nu(x)=\liminf_{r\downarrow0}D_r\nu(x)
\]
（正测度允许取 $+\infty$）。相对于局部有限正 Radon 测度 $\mu$ 的密度只沿中心球
$B(x,r)$、$r\downarrow0$ 观察；只有在 $\mu(B(x,r))>0$ 时才写
$\nu(B(x,r))/\mu(B(x,r))$，并以 $\mu$-a.e. 说明极限。
先对有限 Radon 测度定义
\[
M\nu(x)=\sup_{r>0}\frac{|\nu|(B(x,r))}{|B(x,r)|},
\]
由参数积分的可测性以及从下方逼近半径的有理数列说明 $M\nu$ 可测，再证明与函数情形相同的弱型估计。若
$\nu=g\,m+\nu_s$ 是相对于 Lebesgue 测度的分解（$g\in L^1$、
$\nu_s\perp m$），则完整证明
\[
\lim_{r\downarrow0}\frac{\nu(B(x,r))}{|B(x,r)|}=g(x),
\qquad
\lim_{r\downarrow0}\frac{|\nu_s|(B(x,r))}{|B(x,r)|}=0
\]
对 Lebesgue 几乎处处的 $x$ 成立。局部有限情形用 $B(0,R+1)$ 的限制与可数穷竭推出。

随后建立有界离心率的测度版：若可测集 $V_k\ni x$、$V_k\subset B(x,r_k)$、
$m(V_k)\ge c|B(x,r_k)|$（$0<c\le1$）且 $r_k\downarrow0$，则
\[
\frac{\nu(V_k)}{m(V_k)}\longrightarrow g(x)
\]
对几乎处处 $x$ 成立。取 $h_k\downarrow0$，令
$I_k=[x,x+h_k]$ 与 $[x-h_k,x]$，得到 4.2 所需的单侧增量结论。

相对密度采用一个闭合版本：先相对于 $\mu$ 写
$\nu=\nu_{\rm ac}+\nu_\perp$，其中 $\nu_{\rm ac}=h\mu$、
$\nu_\perp\perp\mu$，并以 $h=\dd\nu_{\rm ac}/\dd\mu$ 记密度。
先给出并证明 Besicovitch 有限重叠选择引理。给定 $R,r_0>0$，称
$\mathcal V$ 是有界集 $E\subset B(0,R)$ 上的细球族，如果每个
$x\in E$ 都是 $\mathcal V$ 中球的中心，并且对每个 $\delta>0$ 存在
$B(x,r)\in\mathcal V$ 满足 $0<r<\min\{\delta,r_0\}$。存在只依赖
维数的 $N(n)$，使得可以从 $\mathcal V$ 选出至多可数的球 $\{B_j\}$，满足

\[
E\subset\bigcup_jB_j,
\qquad
\sum_j\mathbf1_{B_j}\le N(n)\quad\text{处处};
\]

当 $E$ 紧时可选有限子族。正文先实施 Besicovitch 的有限方向算法：在有限截断中按半径递减选取球并删除已覆盖中心，
再把相对于任意测试点的选球按有限个锥方向分组；方向分离与半径比较的打包估计直接给出全族
重叠次数至多为只依赖 $n$ 的 $N(n)$。对一般细球族先用第二可数基作可数编号，再运行同一
有限方向递归；每个有限阶段的覆盖关系与重叠界均采用同一个 $N(n)$，取递归极限后逐点保留
覆盖和重叠界。证明同时记录
所选球仍属于原细球族、覆盖的是 $E$ 而非其闭包，以及有限子族只在紧集情形使用。
质量估计在固定的 $B(0,R)$ 上应用该引理；全空间的结论随后对
$R=1,2,\ldots$ 分别取坏集并取可数并，异常集仍为零测集，各个 $R$ 的选择和估计
分别完成后再合并异常集。
在相对版本中，对有限有符号测度 $\sigma$ 和固定 $r_0>0$ 先定义局部极大函数
\[
M_{\mu,r_0}\sigma(x)=\sup_{\substack{0<r<r_0\\\mu(B(x,r))>0}}
\frac{|\sigma|(B(x,r))}{\mu(B(x,r))}.
\]
分子、分母对中心的参数可测性与有理半径逼近同时给出该上确界的可测代表。
同一有限重叠选择给出相对弱型估计
\[
\mu\{M_{\mu,r_0}\sigma>\lambda\}\le
\frac{N(n)}{\lambda}\,|\sigma|(\mathbb R^n),\qquad\lambda>0,
\]
先在有界球中的紧坏集上证明，再以内正则性撤去紧性；$r_0$ 只控制用于微分的
小尺度。对 $\nu_{\rm ac}=h\mu$，先由第二章 Radon 正则性和第一章的距离函数分离
证明 $C_c(\mathbb R^n)$ 在 $L^1(\mu)$ 中稠密；
连续函数的球比在 $\mu$-几乎处处（$\mu(B(x,r))>0$ 的半径上）由一致连续性趋于
其点值。用上述相对弱型估计控制 $h$ 与连续近似的差异，得到
\[
\frac{(h\mu)(B(x,r))}{\mu(B(x,r))}\longrightarrow h(x)
\quad\text{对 }\mu\text{-a.e. }x.
\]
对 $\nu_\perp\perp\mu$，取 Borel 集 $N$ 使 $\mu(N)=0$，在有界球内用紧集
$K\subset N$ 逼近 $|\nu_\perp|$ 的质量；$x\notin K$ 时小球最终避开 $K$，
剩余质量由相对弱型估计控制，先令半径趋零再令遗漏质量趋零，遂得奇异部分的
相对密度为零。于是，对有限正 Radon 测度 $\mu$ 与有限有符号 Radon 测度 $\nu$，由此得到
\[
\lim_{r\downarrow0}\frac{\nu(B(x,r))}{\mu(B(x,r))}
=\frac{\dd\nu_{\rm ac}}{\dd\mu}(x)
\quad\text{对 }\mu\text{-a.e. }x,
\]
其中只在 $\mu(B(x,r))>0$ 的半径取商；所有正则性步骤先限制到有限测度的有界开集，
再以紧集内逼近实施。
\planitem{证明节奏}
Lebesgue 基准的球密度先按 4.1.3 的逼近--极大函数论证处理；相对版本则固定有限正测度 $\mu$，
先用 $M_{\mu,r_0}$ 的弱型估计把 $h\in L^1(\mu)$ 逼近到连续函数，再以连续函数的球比一致收敛
识别 $h\mu$。奇异部分固定 $R$，在 $W=B(0,R+1)$ 中取承载零集 $N$，再取紧集
$K\subset N\cap W$ 使遗漏的 $|\nu_\perp|$ 质量小。对 $B(0,R)\setminus K$ 上的小球，
用遗漏测度的相对极大弱型估计控制坏点，最后先令阈值、再令遗漏质量趋零。每一步都把
Besicovitch 选择引理的重叠常数 $N(n)$ 写入估计。

\planitem{主例}
\begin{itemize}
\item 对 $g(x)\dd x$，球密度比就是 $A_rg(x)$；对 $\delta_0$，原点的比值为
$(\omega_nr^n)^{-1}\to\infty$，而原点以外最终为零。对三分 Cantor 概率测度 $\mu_C$，
在一维中令 $\alpha=\log2/\log3$；若 $x\in C$、$0<r<1$，由柱区间包含与相交数估计逐步得到
\[
\frac12r^\alpha\le\mu_C(B(x,r))\le6r^\alpha,\qquad
\frac14r^{\alpha-1}\le D_r\mu_C(x)\le3r^{\alpha-1}.
\]
故 $D_r\mu_C(x)\to\infty$ 于 $C$ 上，而在 $m$-几乎处处（事实上 $x\notin C$ 的每一点）
最终为零；基本区间 $I$ 上 $\mu_C(I)=|I|^\alpha$。
\item 对 E4 分别读出原子、奇异连续和绝对连续三项；这组局部结论在 4.2.4 转写成累计函数的三种变化。
\end{itemize}

\planitem{输出与练习}
输出测度微分、单侧区间推论和相对密度定理。练习要求处理有符号原子列、验证局部有限测度的穷竭兼容性、用密度结论证明互相奇异测度在几乎处处的相对密度为零。

\subsection*{4.2　累计链：导数究竟遗漏了什么}

\unithead{4.2.1}{BV、Jordan 分解与端点完整的测度编码}

\planitem{驱动问题}
有限阶梯函数在跳点外导数为零，却能有非零总增量；Cantor 函数导数几乎处处为零，总增量仍为一。需要一个同时记录方向、抵消和奇异质量的累计对象。

\planitem{定义与正式结论}
对 $a<b$ 及实函数 $F:[a,b]\to\R$ 定义
\[
V_a^b(F)=\sup_{a=x_0<\cdots<x_N=b}
\sum_{j=1}^N|F(x_j)-F(x_{j-1})|,
\qquad
\|F\|_{BV}=|F(a)|+V_a^b(F).
\]
记 $BV[a,b]=\{F:[a,b]\to\R:V_a^b(F)<\infty\}$；称 $F$ 右连续，是指每个
$x\in[a,b)$ 都满足 $F(x)=\lim_{y\downarrow x}F(y)$。
先证明区间可加性、变化函数 $v_F(x)=V_a^x(F)$ 的单调性，以及
\[
F=F(a)+p-q,\qquad p=\frac{v_F+F-F(a)}2,
\quad q=\frac{v_F-F+F(a)}2
\]
的函数 Jordan 分解。
当 $F(a)=0$ 时，证明这两个分量在所有以零为端点的非减分解中逐点最小；若 $F$ 取规范右连续代表，
再证明右连续性传给 $p,q$ 及 $v_F$，从而它们可以分别生成 Lebesgue--Stieltjes 测度。
同时证明单调函数在每个内点具有左右极限、跳跃点至多可数且函数 Borel 可测
（沿用命题标签 \texttt{proposition:4-2-1-1}）；这些性质为右连续规范和可数原子列提供代表性依据。

紧接着验证 BV 的代数性质（沿用命题标签
\texttt{proposition:4-2-1-3}）：若 $F,G\in BV[a,b]$、$\alpha\in\mathbb R$，则
$\alpha F,F+G\in BV[a,b]$，且
\[
 V_a^b(\alpha F)=|\alpha|V_a^b(F),\qquad
 V_a^b(F+G)\le V_a^b(F)+V_a^b(G).
\]
由于 $\|F\|_\infty\le |F(a)|+V_a^b(F)$，还可逐分割证明
$FG\in BV[a,b]$ 以及
\[
 V_a^b(FG)\le \|F\|_\infty V_a^b(G)+\|G\|_\infty V_a^b(F).
\]
这些估计先由每个分割上的三角不等式得到，随后承接到 Stieltjes 积分中的乘积估计与分部公式。

随后固定贯穿 4.2 的端点规范。若 $\nu$ 是 $[a,b]$ 上有限有符号 Borel 测度，令
\[
c_\nu=\nu(\{a\}),\qquad F_\nu(a)=0,
\qquad F_\nu(x)=\nu((a,x])\quad(a<x\le b).
\]
反之，给定 $c\in\R$ 与右连续的 $F\in BV[a,b]$、$F(a)=0$，存在唯一测度
\[
\nu=c\delta_a+\nu_F,\qquad
\nu_F((s,t])=F(t)-F(s).
\]
正文证明两个构造互逆，并证明等距关系
\[
|\nu|([a,b])=|c|+V_a^b(F),
\qquad
|\nu_F|((s,t])=V_s^t(F).
\]
记 $F(x+)=\lim_{y\downarrow x}F(y)$（$a\le x<b$）、$F(x-)=\lim_{y\uparrow x}F(y)$（$a<x\le b$；极限由单调 Jordan 分量存在）。
对任意 BV 代表 $F$，令 $G(x)=F(x+)$（$a\le x<b$）且 $G(b)=F(b)$，再以
$c=G(a)$、$\widehat F=G-G(a)$ 得到规范端点二元组。原函数在端点的取值偏差另行列入点态账本。紧接着证明“取值规范化引理”：规范化只在至多可数的不连续点改变取值。分别定义左右极限之间的跳幅以及原值相对右极限的校正幅，证明两族幅度的绝对值总和受总变差控制，再用长度估计说明规范化前后在几乎处处具有相同普通导数。点态变差与测度变差的等距关系只对规范代表陈述；孤立尖点 $\mathbf1_{\{c\}}$ 的左右极限相等、点态变差为 $2$，规范代表为零，专门用来显示跳跃与可去尖点以及两种变差对象的区别。

\planitem{证明节奏}
先在有限阶梯函数上逐分割计算；再对非减函数调用第二章的 Lebesgue--Stieltjes 构造；随后对 Jordan 两部分取测度差，并用 Hahn--Jordan 唯一性证明全变差等式。左端质量始终作为 $c\delta_a$ 显式保留，半开增量只负责 $(a,b]$。

\planitem{主例}
\begin{itemize}
\item 若 $a<x_1<\cdots<x_N<b$ 且
$F=\sum_{k=1}^Nc_k\mathbf1_{[x_k,b]}$，则
\[
\nu_F=\sum_{k=1}^Nc_k\delta_{x_k},
\qquad V_a^b(F)=\sum_{k=1}^N|c_k|.
\]
只要某个 $c_k\ne0$，在跨过该跳点的区间上
$F(x)-F(a)=\int_a^xF'$ 失败；判据由各跳高是否为零决定，净端点增量的抵消与总变差分开记录。
\item 对 $F(x)=x^2\sin(x^{-2})$，$F(0)=0$，取
$x_k=((k+\tfrac12)\pi)^{-1/2}$。则
$F(x_k)=(-1)^k/((k+\tfrac12)\pi)$；按这些点的递减顺序取分割，
相邻增量的绝对值和含有一列调和级数，故给出无限变差；
并写出
\[
F'(x)=2x\sin(x^{-2})-2x^{-1}\cos(x^{-2}),\qquad
\int_0^\eps\frac{2|\cos(x^{-2})|}{x}\,\dd x
=\int_{\eps^{-2}}^\infty\frac{|\cos t|}{t}\,\dd t=\infty,
\]
从而明确 $F'\notin L^1(0,1)$。
\item 将 AC 不定积分、有限阶梯函数和 Cantor 函数排成同一张表，集中计算其变差测度类型；下一单元再由测度微分逐项读出普通导数。
\end{itemize}

\planitem{输出与练习}
输出 BV 的向量与乘积封闭、变差估计、规范代表和端点二元组 $(c,F)$。练习要求合并重复跳点后重算变差、构造符号抵消但变差不抵消的阶梯函数、证明绝对连续/原子/奇异连续三类互相奇异测度的全变差逐项相加。

\unithead{4.2.2}{单调函数可微与 Stieltjes 测试}

\planitem{驱动问题}
累计函数的普通差商是单侧区间上的测度密度。4.1.4 已经控制这种偏心区间，因此可以准确回答：普通导数读到了 Stieltjes 测度的哪一部分？连续测试函数又能读到多少？

\planitem{单调与 BV 的可微性}
若 $F$ 非减，则 $F'$ 几乎处处存在、$F'\ge0$，并且
\[
\int_a^bF'(x)\dd x\le F(b)-F(a).
\]
对规范右连续版本，若 $\nu_F=fm+\nu_s$，则
\[
F'=f\quad\text{a.e.},
\qquad
F(b)-F(a)-\int_a^bF'=\nu_s((a,b]).
\]
对一般 BV 函数分别处理 Jordan 两部分，得到 $F'$ 几乎处处存在、$F'\in L^1$ 且
\[
\int_a^b|F'|\le V_a^b(F).
\]
证明从 4.1.4 的有界离心率测度微分开始：
\[
\frac{F(x+h)-F(x)}h=\frac{\nu_F((x,x+h])}{h},
\qquad
\frac{F(x)-F(x-h)}h=\frac{\nu_F((x-h,x])}{h}.
\]
绝对连续部分给出相同极限，奇异部分的两个商均由中心区间密度的两倍控制。

\planitem{Stieltjes 积分与算子范数}
对连续 $g$ 和规范 $F$ 定义
\[
\int_{(a,b]}g\dd F:=\int_{(a,b]}g\dd\nu_F;
\]
完整测度 $(c,F)$ 的积分为 $cg(a)+\int_{(a,b]}g\dd F$。沿用第一章 \texttt{ex:1-1-1-2} 的带标记分割有向预序，只比较底层分割并允许细化时重选标记；在这一有向集上证明
\[
\sum_{j=1}^Ng(\xi_j)\bigl(F(x_j)-F(x_{j-1})\bigr)
\longrightarrow \int_{(a,b]}g\dd F.
\]
阶梯近似在左端规定为 $g(a)$，从而一致误差确为 $\omega_g(|P|)$。

标量结论之后给出 Banach 值推论。若 $X$ 为 Banach 空间、$g:[a,b]\to X$ 连续，写
$\sigma=\dd\nu_F/\dd|\nu_F|$，于是 $|\sigma|=1$ 于 $|\nu_F|$-几乎处处，并定义 Bochner--Stieltjes 积分
\[
\int_{(a,b]}g\dd F
:=\int_{(a,b]}g(x)\sigma(x)\dd|\nu_F|(x).
\]
由第三章的 Bochner 积分与连续像的紧性证明带标记分割和在 $X$ 的范数中收敛，并得到
\[
\left\|\int_{(a,b]}g\dd F\right\|
\le\int_{(a,b]}\|g\|\dd|\nu_F|.
\]
完整端点对象的积分为 $cg(a)+\int_{(a,b]}g\dd F$，相应估计再加
$|c|\|g(a)\|$。

对 $L_{c,F}(g)=cg(a)+\int g\dd F$ 证明
\[
\|L_{c,F}\|_{C([a,b])^*}=|c|+V_a^b(F).
\]
上界由总变差积分估计；下界对实测度取 Hahn 正负集，以 Radon 正则性选不交紧集 $K_+,K_-$，再用
\[
\varphi(x)=\frac{d(x,K_-)-d(x,K_+)}{d(x,K_-)+d(x,K_+)}
\]
构造 $\|\varphi\|_\infty\le1$ 的连续近似符号函数，最后令遗漏质量趋零。

\planitem{分部积分}
对右连续 $F,G\in BV[a,b]$、$F(a)=G(a)=0$，沿用上述左极限记号，证明精确公式
\[
F(b)G(b)=
\int_{(a,b]}F(x-)\dd\nu_G(x)
+\int_{(a,b]}G(x)\dd\nu_F(x).
\]
证明先做有限阶梯函数，逐个核对共同跳点；一般右连续 BV 函数先把跳点按总跳幅截断，
再对无跳部分作保持统一变差界的折线逼近，并用变差测度的正则性和连续测试函数的统一连续性
传递积分恒等式；尾部原子由其总变差控制。等价的“双右连续积分”版本把
$\sum_x\Delta F(x)\Delta G(x)$ 显式列出。

\planitem{主例}
\begin{itemize}
\item Cantor 函数满足 $C'=0$ a.e.，而 $C(1)-C(0)=1$；差额正是奇异测度质量。
\item 对 $q_k\in(a,b]$ 两两不同、$c_k\ge0$ 且 $\sum_kc_k<\infty$ 的非负纯跳测度
$\sum_kc_k\delta_{q_k}$，由测度微分而非逐项求导得到累计函数导数 a.e. 为零。在满足 $c_k>0$ 的原子点处，原子项给出左差商
$[F(q_k)-F(q_k-h)]/h\ge c_k/h\to+\infty$。
\item 在 $[0,1]$ 上
\[
\int_0^1x^2\dd\bigl(\mathbf1_{[1/2,1]}+x^2\bigr)
=\frac14+\int_0^12x^3\dd x=\frac34.
\]
\item 取 $\nu=\delta_{1/3}-2\delta_{2/3}$，则
$L_\nu(g)=g(1/3)-2g(2/3)$。构造分段线性的连续函数 $\varphi$，使
\[
\|\varphi\|_\infty=1,\qquad
\varphi(1/3)=1,\qquad \varphi(2/3)=-1.
\]
于是 $L_\nu(\varphi)=3=\|\nu\|_{TV}$，直接得到
$\|L_\nu\|=3$，并展示连续测试函数怎样同时分离两个异号原子。
\item 在 $[-1,1]$ 上令 $F=G=\mathbf1_{[0,\infty)}$（此处 $F(0)=G(0)=1$，与章首取 $H(0)=0$ 的 Heaviside 约定区分）。增量测度为 $\delta_0$；
在任意足够细且跨过 $0$ 的分割上，若标记取在 $0$ 左侧，Stieltjes 和的共同跳跃项为 $0$，
若标记恰为 $0$ 或严格取在 $0$ 右侧则为 $1$。两类标记都共尾，却相差 $1$，所以标记分割网不收敛；
这一步把连续被积函数的假设与共同跳跃修正的必要性具体化。
\item 令 $K_m$ 为三分 Cantor 集的第 $m$ 阶柱集，并令
\[
\dd\mu_m=(3/2)^m\mathbf1_{K_m}\dd x.
\]
每个长度 $3^{-m}$ 的基本区间对 $\mu_m$ 与 Cantor 测度 $\mu_C$ 都有质量 $2^{-m}$，故对 $g\in C[0,1]$ 逐区间比较得到
\[
\left|\int g\dd\mu_m-\int g\dd\mu_C\right|
\le \omega_g(3^{-m})\longrightarrow0.
\]
这一计算把绝对连续近似的 Stieltjes 积分送到奇异连续极限，并与原子例组成三类测度的完整对照。
\end{itemize}

\planitem{输出与练习}
输出 $F'=\dd\nu_{F,\rm ac}/\dd m$、标量及 Banach 值 Stieltjes 网和、范数等式和完整跳跃修正。练习要求由端点二元组计算含左端原子的积分、为三个原子构造范数达到函数、估计 Cantor 柱状近似的连续模误差，并把分部积分公式化为 AC 函数的熟悉形式。

\unithead{4.2.3}{绝对连续、微积分基本定理与复合边界}

\planitem{驱动问题}
第三章已经把绝对连续函数表示为不定积分；本单元用微分定理识别其中的密度，并回答什么时候“积分后求导”和“求导后积分”真正互逆。

沿用第三章的定义：$F:[a,b]\to\R$ 称为绝对连续（记 $F\in AC[a,b]$），若对每个
$\varepsilon>0$ 存在 $\delta>0$，使任意有限个两两不交的开区间
$(\alpha_j,\beta_j)\subset[a,b]$ 在
$\sum_j(\beta_j-\alpha_j)<\delta$ 时都满足
$\sum_j|F(\beta_j)-F(\alpha_j)|<\varepsilon$。
先由这一 $\varepsilon$--$\delta$ 条件证明短区间族的总变差可作同样控制，因而
$AC[a,b]\subset BV[a,b]$；随后调用第三章的积分表示。
称函数 $F$ 具有 Luzin $N$ 性质，是指每个 Lebesgue 零测集
$E\subset[a,b]$ 都满足 $m^*(F(E))=0$；这一性质将在本单元证明，并在
4.2.4 用来刻画连续 BV 函数的绝对连续性。

\planitem{正式结论}
引用第三章已经证明的表示
\[
F\in AC[a,b]
\quad\Longleftrightarrow\quad
F(x)=F(a)+\int_a^xh(t)\dd t\quad(h\in L^1(a,b)).
\]
由 4.1.3 得 $F'=h$ a.e.，从而闭合两条微积分基本定理：
\[
\frac{\dd}{\dd x}\int_a^xh(t)\dd t=h(x)\quad\text{a.e.},
\qquad
F(y)-F(x)=\int_x^yF'(t)\dd t,
\]
并证明
\[
V_a^b(F)=\int_a^b|F'(t)|\dd t.
\]
在此基础上依次证明。$N$ 性质的证明取零测集的开覆盖使区间总长度任意小；
对每个覆盖区间作有限细分，以连续性把其像集覆盖为增量长度之和，再对有限子族应用
绝对连续条件并取单调极限，得到像集外测度任意小：
\begin{enumerate}
\item AC 函数一致连续、具有 Luzin $N$ 性质，并对加法、数乘和区间限制封闭；
Lipschitz 函数是其中的一个定量子类，且把零测集映成零测集；$\sqrt{x}$ 给出 AC
严格大于 Lipschitz 的可算例。
\item 若 $G\in AC[a,b]$ 且外函数 $\Phi$ 在包含 $G([a,b])$ 的区间上为 $C^1$，则
$\Phi\circ G\in AC[a,b]$，并在 $G$ 可微处成立
$(\Phi\circ G)'=\Phi'(G)G'$；绝对连续性由 $\Phi$ 在紧像集上的 Lipschitz 界和
$G$ 的短区间控制给出。
\item 对任意 $G\in AC[a,b]$ 与包含 $G([a,b])$ 的区间上的连续函数 $h$，令
 $H(y)=\int_{G(a)}^yh(s)\dd s$。先由 $H$ 的 Lipschitz 性和 AC 的短区间控制证明
 $H\circ G\in AC$，再在可微点用链式法则并调用 FTC，得到带方向区间的换元公式
\[
\int_a^bh(G(x))G'(x)\dd x
=\int_{G(a)}^{G(b)}h(y)\dd y.
\]
 右端按有向区间理解；单调非减且 AC 情形的非负 Borel 被积函数版本留待
 4.2.4 的面积公式之后由单调收敛推出。
\item 先定义向量曲线的绝对连续性：$\gamma:[a,b]\to\R^d$ 满足对任意
$\varepsilon>0$ 存在 $\delta>0$，使任意有限个两两不交的区间族满足
$\sum_j(\beta_j-\alpha_j)<\delta$ 时有
$\sum_j|\gamma(\beta_j)-\gamma(\alpha_j)|<\varepsilon$；这等价于各分量绝对连续，
并约定 $\gamma'$ 为各分量导数组成的向量。对这样的 AC 曲线定义
$\operatorname{Len}(\gamma)=\sup_P\sum_j|\gamma(t_j)-\gamma(t_{j-1})|$，再证明
$\operatorname{Len}(\gamma)=\int_a^b|\gamma'|$；证明从多边形长度上界和简单函数近似的反向估计两边夹逼。
\end{enumerate}

\planitem{换元算例}
取 $G(t)=t^2$、$h(s)=(1+s)^{-1}$ 于 $[0,1]$，由带方向换元公式并直接积分核对
\[
\int_0^1\frac{2t}{1+t^2}\dd t
=\int_0^1\frac{\dd s}{1+s}=\log2.
\]
这个有限计算同时检验导数因子、端点方向和连续被积函数的条件。

\planitem{复合反例}
在 $I_k=(2^{-k-1},2^{-k})$ 上作端点取零、中点高度 $a_k=k^{-2}$ 的三角帐篷 $\psi_k$，令 $G=\sum_k\psi_k$、$G(0)=0$。支撑不交、$a_k\to0$ 保证零点连续；有限首段逐段 Lipschitz，尾段总变差 $\sum_{k\ge K}2a_k\to0$，所以按绝对连续定义 $G\in AC$，并且
\[
\int_0^1|G'|=\sum_k2a_k<\infty,
\]
故 $G$ 绝对连续。外函数 $\Phi(s)=\sqrt{s}$ 也绝对连续，但
\[
V_0^1(\Phi\circ G)\ge\sum_k2\sqrt{a_k}
=2\sum_k\frac1k=\infty.
\]
这个例子精确说明“两函数均 AC”本身不足以保证复合仍 AC，并反向解释外层 Lipschitz 假设的作用。

\planitem{输出与练习}
输出 FTC 的双向闭合、变差等式、换元和曲线长度。练习要求验证 $x^\alpha$ 的 AC/Lipschitz 阈值、构造有限个帐篷时的变差公式、在单调换元中处理平坦区间。

\unithead{4.2.4}{三类变化、indicatrix、面积公式与 Banach--Zarecki}

\planitem{驱动问题}
普通导数只读出绝对连续密度。怎样把跳跃和奇异连续变化分开？对一个连续 BV 函数，什么可检验条件恰好排除 Cantor 型遗漏？

\planitem{Luzin $N$ 性质的回访}
沿用 4.2.3 的定义：对连续函数 $F:[a,b]\to\mathbb R$，每个 Lebesgue 零测集
$E\subset[a,b]$ 都满足 $m^*(F(E))=0$。这一性质只涉及像集的测度，
将在本单元作为区分绝对连续变化与奇异连续变化的可检验条件。

\planitem{三分解定理}
对右连续、$F(a)=0$ 的实 BV 函数，将增量测度唯一分解为
\[
\nu_F=fm+\sum_{x\in(a,b]}a_x\delta_x+\nu_{sc},
\]
其中 $f\in L^1$、$\sum_x|a_x|<\infty$，$\nu_{sc}\perp m$ 且无原子。相应地
\[
F(x)=\int_a^xf(t)\dd t
+\sum_{a<y\le x}a_y+\nu_{sc}((a,x]),
\]
三项分别为绝对连续、纯跳跃和奇异连续部分，且 $F'=f$ a.e.。证明完全由第二章测度分解、原子至多可数和 4.2.1 的端点对应组成。

\planitem{Banach indicatrix}
本单元中凡出现连续 BV 函数 $F$，先约定
$\nu_F:=\nu_{F-F(a)}$，即 $\nu_F((s,t])=F(t)-F(s)$；因连续性，
$|\nu_F|$ 在端点单点上无质量。
对连续 BV 函数定义层集重数
\[
N(F,E,y)=\#\{x\in E:F(x)=y\},
\]
端点的计数差异只影响有限个高度；局部公式中的 $V_{\alpha}^{\beta}(F)$ 解释为
$|\nu_F|((\alpha,\beta))$，对连续 $F$ 与通常闭区间变差相同。正文先建立并证明一维计数可测性引理：
对紧 $E$，各超水平集由
$E^k$ 中彼此分离的有限点组的紧像组成；对一般 Borel $E$，相同的点组条件给出
Borel 集在高度轴上的投影，先用紧集内逼近，再用“紧区间中 Borel 投影是
Lebesgue 可测”的投影引理（正文给出紧核逼近与外正则性证明）得到可测性。
因此正文证明 $y\mapsto N(F,E,y)$ 的 Lebesgue 可测性，并以此作为后续积分的可测基础。
随后按折线函数、有限开区间并、开集、Borel 集四步建立局部 indicatrix 公式：若
$U=\bigcup_{i=1}^m(\alpha_i,\beta_i)$ 是有限个两两不交的相对开区间并，则
\[
|\nu_F|(U)=\sum_{i=1}^mV_{\alpha_i}^{\beta_i}(F)
 =\int_\R N(F,U,y)\dd y.
\]
再以正则性和单调类延至任意 Borel 集；全区间版本为
\[
V_a^b(F)=\int_\R N(F,[a,b],y)\dd y.
\]

\planitem{测度延拓引理}
在把局部公式延至 Borel 集以前，先陈述并证明有限测度的
$\pi$--$\lambda$ 唯一性：若 $\mathcal P$ 含全集、对有限交封闭并生成 $\sigma$-代数
$\mathcal A$，且两个有限非负测度在 $\mathcal P$ 上相等，则它们在 $\mathcal A$
上相等。证明把两测度取值相等的集合组成的类验证为含全集、对补集和可数不交并
封闭的 $\lambda$-系统，再用由 $\mathcal P$ 生成的最小 $\lambda$-系统等于
$\sigma(\mathcal P)$；同一验证也给出单调类版本。后文所有“正则性延拓”和面积
公式的测度唯一性均明确调用这一引理。

\planitem{一维面积公式与复合链式法则}
对 Borel 集 $E\subset[a,b]$ 定义同一重数函数 $N(G,E,y)$，取值于
$\{0,1,2,\ldots,\infty\}$。先调用上一段的一维计数可测性引理，再在紧 $E$ 上用
紧集的有限点组像、对开集作可数区间并，并以 Borel 正则性和单调类延至任意 Borel
$E$。本单元的
面积公式均先在 Borel 集上陈述；对 AC 函数若需完成测度，只在确认其 $N$ 性质后取
Borel 代表。对任意连续 BV 函数 $G$ 沿用上述约定，记
$\nu_G:=\nu_{G-G(a)}$，亦即 $\nu_G((s,t])=G(t)-G(s)$；这是
4.2.1 的端点规范在一般基值下的简写。正文由折线函数的有限单调段换元、分割的水平线
计数上下估计和变差测度的正则性，逐步证明连续 BV 函数的 indicatrix 公式，并在 AC
情形识别其变差测度的密度，从而得到面积公式
\[
|\nu_G|(E)=\int_\R N(G,E,y)\,\dd y\quad(G\text{ 连续 BV}),
\qquad
|\nu_G|(E)=\int_E|G'(x)|\,\dd x\quad(G\in AC[a,b]).
\]
其中第二个等式由区间变差等式和测度唯一性先得到
$|\nu_G|=|G'|m$ 于所有 Borel 集。两边的测度结构在正文中逐项写出：
先在紧区间和有限开区间并上核对折线及连续 BV 的等式；对两两不交的
$E_i$，逐个高度有 $N(G,\bigcup_iE_i,y)=\sum_iN(G,E_i,y)$，由单调收敛可知
$E\mapsto\int_\R N(G,E,y)\,\dd y$ 与 $E\mapsto|\nu_G|(E)$ 都是 Borel
$\sigma$-代数上的非负测度。再用区间半环上的一致性和本单元的
$\pi$--$\lambda$ 引理延至所有 Borel 集。对 AC 函数在确认 $N$ 性质后再取
Lebesgue 完成测度；$G'$ 在不可微零测集上置为零，右端按扩展非负积分理解。

取一维 Rademacher 引理为本单元的内部引理：Lipschitz 函数在 a.e. 点可微，且其
导数可取 Borel 代表。给定 $G\in AC[a,b]$ 及在 $G([a,b])$ 邻域上 Lipschitz 的
外函数 $\Phi$，令 $Z$ 为 $\Phi$ 在紧像集 $G([a,b])$ 上的不可微值集的 Borel 零测
超集。将 $\Phi$ 在像集外作 Lipschitz 延拓，并记其 Lipschitz 常数为 $L$。对任意
两两不交的区间族有
\[
\sum_j|\Phi(G(\beta_j))-\Phi(G(\alpha_j))|
\le L\sum_j|G(\beta_j)-G(\alpha_j)|
\]
故 AC 的 $\varepsilon$--$\delta$ 条件先给出 $\Phi\circ G\in AC$；随后面积公式用于
$E=G^{-1}(Z)$ 给出
\[
\int_{G^{-1}(Z)}|G'|\,\dd x=0,
\]
故 $G'=0$ 在该层上 a.e.；在其余 $G$ 可微点用差商得到：若 $\Phi$ Lipschitz，
则 $\Phi\circ G\in AC$ 且
\[
(\Phi\circ G)'=\Phi'(G)G'\quad\text{a.e.},
\]
并把 $\Phi'$ 在 $Z$ 上置为零。对单调非减且绝对连续的 $G$，面积公式再以平坦层的零导数
贡献和单调收敛推出任意非负 Borel $h$ 的值域换元公式；连续单调函数的换元与
4.2.3 的连续 $h$ 公式在同一端点方向约定下衔接。若 $G$ 为右连续的单调非减 BV 函数，
令 $J_G=\{x\in(a,b]:\Delta G(x)>0\}$（至多可数），并按 4.2.2 的 Stieltjes 测度定义
$\int h(G)\,\dd G$。对有界非负 Borel $h$ 写出完整的跳跃修正，并对一般非负 $h$
通过截断和单调收敛延拓：
\[
\int_{(a,b]}h(G(x))\,\dd G(x)
=\int_{G(a)}^{G(b)}h(y)\,\dd y
+\sum_{x\in J_G}\left[h(G(x))\Delta G(x)-
\int_{G(x-)}^{G(x)}h(y)\,\dd y\right].
\]
第一项减去各跳跃区间的积分后正是连续部分的值域积分，方括号中的第一项是每个
Stieltjes 原子的贡献，从而跳跃与连续增长在公式中分别可见；对一般非负 $h$，右端按
$h_N=h\wedge N$ 所得整个表达式的单调极限理解。

局部 indicatrix 还用于证明并正式编号“变差--$N$ 引理”：若连续 BV 函数 $F$
把零测集映成零测集，则 $v_F(x)=V_a^x(F)$ 连续且也具有该性质。先由连续性说明
$\nu_{v_F}=|\nu_F|$ 无原子；若 $E$ 是任意 Lebesgue 零测集，取包含它的 Borel 零测包
$B\subset[a,b]$，则 $F(B)$ 仍为零测集，Borel 版面积公式给出
$|\nu_F|(B)=\int_\R N(F,B,y)\,\dd y=0$，再由单调性得到 $|\nu_F|(E)=0$。对连续非减的
$v_F$，先对 Borel 集 $B$ 用单调函数的区间像覆盖（等价地，其一维面积公式）得到
$m^*(v_F(B))\le\nu_{v_F}(B)=0$；由 $E\subset B$ 遂有 $m^*(v_F(E))=0$。
这一引理随后作为 Banach--Zarecki 充分性的直接输入。

随后完整证明 Banach--Zarecki 定理：
\[
F\in AC[a,b]
\quad\Longleftrightarrow\quad
F\text{ 连续、属于 }BV[a,b]\text{ 且具有 Luzin }N\text{ 性质}.
\]
必要性用 AC 的小区间定义；充分性先由继承引理得到 $v_F$ 具有 $N$ 性质，再证明连续非减且有 $N$ 性质的函数其 Stieltjes 测度绝对连续，因此 $v_F\in AC$。由
$|F(y)-F(x)|\le v_F(y)-v_F(x)$ 把 $v_F$ 的绝对连续性传给 $F$；等价地，Jordan 两部分的增量均受 $v_F$ 增量控制。最后调用第三章的 AC 积分表示识别导数。

\planitem{主例}
\begin{itemize}
\item E4 的规范累计函数具有非零的三类成分；计算其 a.e. 导数恒为 $1-p-q$，总变差则读取三项测度的全部质量。
\item Cantor 函数把零测 Cantor 集映到 $[0,1]$，因此恰在 $N$ 性质上失败。
\item 在 $I_k=(2^{-k-1},2^{-k})$ 上令 $\psi_k$ 为端点取零、中点高度为 $2^{-k}$ 的线性帐篷，并置 $F=\sum_{k\ge1}\psi_k$。支撑不交且在零点连续，直接计算
\[
V_0^1(F)=2\sum_{k\ge1}2^{-k}=2,
\qquad N(F,[0,1],0)=\infty.
\]
对 $y>0$ 有
$N(F,[0,1],y)=2\#\{k:2^{-k}>y\}$（峰值高度组成的可数例外不影响积分），逐层积分重新得到 $2$，从而以完整计算检验 indicatrix 公式及“无穷重数集中在零测高度集”的结论。
\end{itemize}

\planitem{本节闭环与后续接口}
4.2 的终点是“累计函数 $\leftrightarrow$ 端点原子加 Stieltjes 测度 $\leftrightarrow$ 连续测试积分”。第七章在整轴上的概率分布语境中，以这组编码再加入对无穷远质量的统一控制，建立 Helly 选择和分布函数收敛判据。本章末练习聚焦紧区间上的直接推论：若 $(c_k,F_k)$ 的总变差一致有界且连续测试积分已经收敛，则极限泛函的范数满足下半连续；证明使用 4.2.2 的范数公式。

\planitem{输出与练习}
输出 BV 三分解、indicatrix、AC 面积公式、Lipschitz 链式法则和 Banach--Zarecki。练习要求给有限折线逐层计算重数积分、验证平坦段只贡献零测高度、分别改变 E4 的三个参数并判断普通导数与总变差如何变化。

\subsection*{4.3　测试链：从卷积光滑化到 Sobolev 紧性}

\unithead{4.3.1}{Young 不等式、mollifier 与光滑稠密}

\planitem{驱动问题}
第三章已经证明 $L^1$ 卷积几乎处处存在及其 $L^1$ 估计。本单元从这个接口出发，考察平移平均在什么指数范围内稳定、怎样制造光滑近似，以及点态边界为何与范数恢复不同。

\planitem{卷积估计}
先回引第三章的 Borel 代表、Tonelli--Fubini、$L^p$ 平移连续性和 $L^1*L^1$ 结论。正文证明基础 Young 估计
\[
\|f*g\|_p\le\|f\|_1\|g\|_p,
\qquad f\in L^1,\quad g\in L^p,\quad1\le p\le\infty,
\]
以及一般 Young 不等式：若 $1\le p,q,r\le\infty$ 且
\[
1+\frac1r=\frac1p+\frac1q,
\]
则
\[
\|f*g\|_r\le\|f\|_p\|g\|_q.
\]
证明把两个端点、$r=\infty$ 的 H\"older 情形和中间指数分开，再用加权 H\"older 与 Fubini 完成；每次交换积分前标出非负性或绝对可积上界。当 $r=\infty$ 时，平移具有有限指数的那个因子，并以另一因子的范数控制差值，得到有界连续代表；端点 $(p,q)=(1,\infty)$ 与 $(\infty,1)$ 也按这一方式处理。

\planitem{mollifier 与恢复}
取 $\rho\in C_c^\infty(\R^n)$、$\rho\ge0$、$\int\rho=1$，令
$\rho_\eps(x)=\eps^{-n}\rho(x/\eps)$。证明：
\begin{quote}
记 $\mathbb N_0=\{0,1,2,\ldots\}$。多重指标先约定为
$\alpha=(\alpha_1,\ldots,\alpha_n)\in\mathbb N_0^n$，
$|\alpha|=\sum_j\alpha_j$，$\partial^\alpha=\partial_1^{\alpha_1}\cdots
\partial_n^{\alpha_n}$，并把 $\partial^0$ 记作恒等算子。
\end{quote}
\begin{enumerate}
\item $f\in L^1_{\rm loc}$ 时，$f*\rho_\eps$ 在允许的内部区域光滑，且
\[
\partial^\alpha(f*\rho_\eps)=f*(\partial^\alpha\rho_\eps).
\]
\item $1\le p<\infty$ 时
$\|f*\rho_\eps-f\|_p\to0$；$f\in C_0$ 时一致收敛。范数证明写成
\[
\|f*\rho_\eps-f\|_p
\le\int\rho_\eps(y)\|\tau_yf-f\|_p\dd y
\]
并直接调用第三章平移连续性。
若 $f\in L^p_{\rm loc}(\Omega)$ 且 $K\Subset\Omega$，先以一个在 $K$ 邻域恒为一的截断作局部延拓，再取小于该邻域边距的 $\eps$，得到
$\|f*\rho_\eps-f\|_{L^p(K)}\to0$。
端点 $p=\infty$ 单独核对：对 Heaviside 函数 $H=\mathbf1_{(0,\infty)}$ 和偶标准核，
$\|H*\rho_\eps-H\|_\infty\ge\tfrac12$；因此跳跃函数在 $L^\infty$ 中不发生范数恢复。
\item 若 $x$ 是 $f$ 的 Lebesgue 点，则 $f*\rho_\eps(x)\to f(x)$；证明以支撑球内的平均振荡控制。
\item 对任意开集 $\Omega\subset\R^n$，$C_c^\infty(\Omega)$ 在 $L^p(\Omega)$ 中稠密（$1\le p<\infty$）：先用 Radon 内正则性把质量截在 $K\Subset\Omega$，再取内部截断并以小于 $\dist(K,\Omega^c)$ 的尺度卷积，逐项控制截断误差、卷积误差与支撑。该结论沿用标签 \texttt{theorem:4-3-1-3}。
\end{enumerate}

局部论证统一使用内部耗尽：对 $\eps>0$ 置
\[
 \Omega^{\mathrm{ex}}_\eps
 =\{x\in\Omega:\dist(x,\Omega^c)>\eps,\ |x|<\eps^{-1}\}.
\]
它是开集，$\overline{\Omega^{\mathrm{ex}}_\eps}\Subset\Omega$，且
$\Omega^{\mathrm{ex}}_{1/k}\uparrow\Omega$（沿用定义标签 \texttt{def:4-3-1-3}）。这样每次先在一个固定紧内缩上完成卷积或测试函数论证，
再令 $k\to\infty$，边界与无穷远分别由内缩参数控制。

\planitem{平移接口}
直接调用第三章的平移连续命题
\texttt{prop:3-3-3-translation-continuity}（第三章写作 $\tau_hf(x)=f(x+h)$；
本计划的 $\tau_hf(x)=f(x-h)$ 由替换 $h\mapsto-h$ 得到同一结论），并在本单元把它逐截面提升为下述有限测度乘积空间版本：若
$(S,\Sigma,\nu)$ 为有限测度空间、$1\le p<\infty$、$F\in L^p(S\times\R)$，则
\[
\|F(\cdot,\cdot+h)-F\|_{L^p(S\times\R)}\longrightarrow0.
\]
对几乎处处的截面使用一维平移连续性，以
\[
\|F(s,\cdot+h)-F(s,\cdot)\|_p^p
\le2^{p-1}\bigl(\|F(s,\cdot+h)\|_p^p+\|F(s,\cdot)\|_p^p\bigr)
\]
作共同可积控制。这个推论直接供第六章的时间变量平移调用，并在 4.4.1 作为
\texttt{theorem:4-4-1-1} 统一登记。

\planitem{主例与边界}
\begin{itemize}
\item 取 $f(0)=g(0)=0$，并在 $x\ne0$ 时令
$f(x)=g(x)=|x|^{-1/2}\mathbf1_{(-1,1)}(x)$，则二者属于 $L^1(\R)$，但
$(f*g)(0)=\int_{-1}^1|y|^{-1}\dd y=+\infty$。结论准确落在“几乎处处绝对收敛”。
\item 计算 $\mathbf1_{(0,1)}*\mathbf1_{(0,1)}$ 的三角形图像，所得公式在 4.4 直接用于检验支撑与连续性。
\item 对满足 $\rho(x)=\rho(-x)$ 的对称标准核，零延拓的 $\mathbf1_{(0,1)}$ 在边界点 $0$ 的卷积值为 $1/2$；每个固定内部点最终收敛到 $1$，而 $x_\eps/\eps\to0$ 的移动点仍看见边界过渡层。
\item 当 $1\le p,q<\infty$ 且 $1/p+1/q<1$，选择 $a>n/p$、$b>n/q$ 且 $a+b\le n$，令
$f=(1+|x|)^{-a}$、$g=(1+|x|)^{-b}$；逐项验证因子所属空间，而卷积尾积分发散，说明 Young 指数关系的范围。
\end{itemize}

\planitem{输出与练习}
输出一般 Young、全空间及局部 mollifier 收敛、任意开集上的 $C_c^\infty$ 稠密和乘积空间平移。练习要求证明卷积平滑的导数范数代价
$\|\partial^\alpha(f*\rho_\eps)\|_p\le
\eps^{-|\alpha|}\|\partial^\alpha\rho\|_1\|f\|_p$，比较固定点与移动点的边界极限，并把乘积空间推论推广到 $\sigma$-有限且乘积测度已构造的情形。

\unithead{4.3.2}{分布、分布导数与弱极限}

\planitem{驱动问题}
尖峰的 $L^1$ 极限可能不再是函数，阶跃的经典导数忽略跳跃。若只要求导数对光滑紧支撑测试函数作用，二者能否进入同一个闭合框架？

\planitem{定义次序}
先记 $\mathcal D(\Omega)=C_c^\infty(\Omega)$，$\mathcal D'(\Omega)$ 为其连续对偶，再定义测试函数的局部半范数。
分布 $T\in\mathcal D'(\Omega)$ 是 $\mathcal D(\Omega)$ 上的线性泛函，且对每个 $K\Subset\Omega$，存在 $C_K>0$ 与整数 $m_K$，使
\[
|\langle T,\varphi\rangle|
\le C_K\sum_{|\alpha|\le m_K}\|\partial^\alpha\varphi\|_\infty,
\qquad \supp\varphi\subset K.
\]
随后对 $f\in L^1_{\rm loc}(\Omega)$ 定义正则分布

\[
\langle T_f,\varphi\rangle=\int_\Omega f(x)\varphi(x)\,\dd x.
\]

每个局部有限 Radon 测度 $\mu$ 通过
\[
\langle T_\mu,\varphi\rangle=\int_\Omega\varphi\dd\mu,
\qquad
|\langle T_\mu,\varphi\rangle|
\le |\mu|(K)\|\varphi\|_\infty
\quad(\supp\varphi\subset K\Subset\Omega)
\]
成为零阶分布；Dirac 分布和 4.2 的 Stieltjes 测度 $\nu_F$ 都作为这一嵌入的具体对象出现。
规定 $T_k\to T$ 在分布意义下，当且仅当对每个
$\varphi\in C_c^\infty(\Omega)$ 都有

\[
\langle T_k,\varphi\rangle\longrightarrow\langle T,\varphi\rangle.
\]

对 $1\le p<\infty$，先约定 $u_k\weakto u$ 于 $L^p$ 是指
$\int u_k\phi\dd x\to\int u\phi\dd x$ 对每个 $\phi\in L^{p'}$ 成立
（$p=1$ 时 $p'=\infty$）。本单元先证明并回引弱序列紧性引理：
$1<p<\infty$ 时，$L^p$ 中的有界序列有弱收敛子列；证明取
$L^{p'}$ 的可数稠密族作对角抽取，把极限延拓为有界线性泛函，再由第三章的
对偶表示识别为某个 $L^p$ 函数。

再定义分布导数
\[
\langle\partial_jT,\varphi\rangle=-\langle T,\partial_j\varphi\rangle.
\]
若 $\supp\rho\subset B(0,1)$，置
$\Omega_\eps=\{x\in\Omega:\dist(x,\Omega^c)>\eps\}$，并对任意
$T\in\mathcal D'(\Omega)$ 定义局部光滑化
\[
(T*\rho_\eps)(x)=\langle T,\rho_\eps(x-\cdot)\rangle,
\qquad x\in\Omega_\eps.
\]
光滑函数乘分布也在给出作用公式及支撑条件后使用：对 $a\in C^\infty(\Omega)$ 与测试函数
$\varphi$，$a\varphi$ 仍为紧支撑测试函数且
$\supp(a\varphi)\subset\supp\varphi$。

\planitem{正式结论}
\begin{enumerate}
\item 分布导数永远存在（沿用标签 \texttt{proposition:4-3-2-1}）：对每个多重指标 $\alpha$，$\partial^\alpha T$ 仍为分布，并且
\[
 \partial^\alpha\partial^\beta T=\partial^{\alpha+\beta}T.
\]
证明在固定紧集的半范数估计中把测试函数的一个偏导移入已有估计，再用测试函数的经典混合偏导交换；
因此任意次微分均有定义，求导次序不产生歧义。
\item 对 $a\in C^\infty(\Omega)$ 定义
\[
\langle aT,\varphi\rangle=\langle T,a\varphi\rangle.
\]
证明乘积仍为分布，并核对 Leibniz 公式
\[
\partial_j(aT)=a\,\partial_jT+(\partial_ja)T.
\]
这条公式随后与 Sobolev 中的截断和乘积法则相接。
\item 分布光滑化：$T*\rho_\eps\in C^\infty(\Omega_\eps)$，并且
\[
\partial^\alpha(T*\rho_\eps)
=(\partial^\alpha T)*\rho_\eps
=T*(\partial^\alpha\rho_\eps).
\]
证明把差商作用到测试函数平移上，并在每个固定紧集的半范数中控制余项。对 $T=T_\mu$，直接算得
\[
(T_\mu*\rho_\eps)(x)=\int\rho_\eps(x-y)\dd\mu(y);
\]
4.4.1 沿用这一公式作为函数--测度卷积。特别地，$T_{\nu_F}*\rho_\eps$ 把 4.2 的累计变化送入测试链的光滑对象。
\item 基本引理：若 $f\in L^1_{\rm loc}$ 且 $\int f\varphi=0$ 对所有测试函数成立，则 $f=0$ a.e.；证明用 mollifier 和局部恢复，保证函数型弱导数唯一。
\item 闭性：若 $f_k\to f$、$g_{k,j}\to g_j$ 于 $L^1_{\rm loc}$，且 $\partial_jT_{f_k}=T_{g_{k,j}}$，则 $\partial_jT_f=T_{g_j}$。
\item 对 $1<p<\infty$，令 $p'=p/(p-1)$。若
$\partial_jT_{f_k}=T_{g_{k,j}}$、$T_{f_k}\to T_f$ 于分布意义且
$\sup_k\|g_{k,j}\|_p<\infty$，则以 $L^{p'}$ 的可分稠密族作对角抽取，
调用上述弱序列紧性引理得到共同子列弱收敛于 $g_j\in L^p$，并识别
$\partial_jT_f=T_{g_j}$。
\end{enumerate}
最后一项的端点边界由实际序列说明：在 $(-1,1)$ 上令
$f_k(x)=0$（$x\le0$）、$f_k(x)=kx$（$0<x<1/k$）、$f_k(x)=1$（$x\ge1/k$）。
则 $f_k\to\mathbf1_{(0,\infty)}$ 于 $L^1_{\rm loc}$，而
$f_k'=k\mathbf1_{(0,1/k)}$ 在 $L^1$ 有界并分布收敛到 $\delta_0$；
$p=1$ 时共轭测试空间 $L^\infty$ 不可分，且有界集一般缺少所需的弱相对紧性。

\planitem{主例}
逐次计算
\[
H'=\delta_0,
\qquad (|x|/2)''=\delta_0,
\qquad \partial^\alpha\delta_0(\varphi)=(-1)^{|\alpha|}\partial^\alpha\varphi(0).
\]
每个计算都从测试函数积分分部开始，并核对边界项。$|x|/2$ 的一维基本解例完整承载“奇性由弱导数编码”的动机，同时连接后面的热核生成元。

\planitem{输出与练习}
输出测度的零阶分布嵌入、任意阶分布导数及其混合偏导交换、光滑乘法、任意分布的局部光滑化以及函数型弱导数的唯一性与闭性。练习要求计算分段 $C^1$ 函数的分布导数、证明平移与分布导数交换、判断一列对称尖峰的分布极限。

\unithead{4.3.3}{Sobolev 基础、一维代表、ACL 与 Meyers--Serrin}

\planitem{驱动问题}
若一个分布导数重新由 $L^p$ 函数表示，怎样把它组织成完备函数空间？弱导数在哪些方向重新成为普通导数？

\planitem{定义与完备性}
对开集 $\Omega\subset\R^n$、$1\le p\le\infty$ 定义
\[
W^{1,p}(\Omega)=\{u\in L^p(\Omega):
\partial_ju\text{ 由 }L^p\text{ 函数表示},\ j=1,\ldots,n\},
\]
配以 $\|u\|_{W^{1,p}}=\|u\|_p+\sum_j\|\partial_ju\|_p$；
$W_0^{1,p}(\Omega)$ 定义为 $C_c^\infty(\Omega)$ 的闭包。利用 $L^p$ 完备性与 4.3.2 的闭性逐坐标证明 $W^{1,p}$ 完备。

\planitem{代表定理}
先约定 ACL（绝对连续截面）代表的含义：在几乎每条与坐标轴平行的线段上，
限制函数具有一维绝对连续代表，且各坐标方向的例外线集取一个共同零测集。
\begin{enumerate}
\item 若 $I=(a,b)$ 有限且 $u\in W^{1,p}(I)$，则存在唯一连续 AC 代表
\[
\widetilde u(x)=c+\int_a^xu'(t)\dd t.
\]
证明先由分布导数为零推出函数 a.e. 为常数，再调用第三章 AC 表示。
\item 对长方体 $Q\Subset\Omega$，由 Fubini 对几乎每条坐标线识别一维弱导数，得到 ACL 代表；例外线集、代表选择和各方向共同零测集逐项写出。
\item 乘积法则：若 $a\in C^1$ 且 $a,\nabla a\in L^\infty$，则
$\partial_j(au)=a\partial_ju+u\partial_ja$。先在分布配对中证明，随后估计范数。
\item Sobolev 卷积求导：若 $u\in W^{1,p}(\R^n)$、$g\in L^1(\R^n)$，其中 $1\le p\le\infty$，则
\[
u*g\in W^{1,p}(\R^n),\qquad
\partial_j(u*g)=(\partial_ju)*g,
\qquad
\|\partial_j(u*g)\|_p
\le\|\partial_ju\|_p\|g\|_1.
\]
在测试函数配对中依次使用 Fubini、平移换元和弱导数定义，再由 Young 不等式完成估计。特别地，若 $u,g\in W^{1,1}$，则
\[
\partial_j(u*g)=(\partial_ju)*g=u*(\partial_jg).
\]
\end{enumerate}

\planitem{Meyers--Serrin}
对任意开集和 $1\le p<\infty$ 证明
\[
C^\infty(\Omega)\cap W^{1,p}(\Omega)
\quad\text{在 }W^{1,p}(\Omega)\text{ 中稠密}.
\]
证明先由第一章的光滑 bump 构造与局部有限穷竭相容的光滑单位分解；每一块的支撑紧包含于 $\Omega$ 后作零延拓与小尺度卷积；按 $2^{-j}\eps$ 分配误差，再局部有限求和。紧接着将两种逼近类并列：一般 $W^{1,p}$ 的逼近类是 $C^\infty(\Omega)\cap W^{1,p}(\Omega)$，而 $C_c^\infty(\Omega)$ 的闭包定义 $W_0^{1,p}(\Omega)$。

\planitem{主例}
\begin{itemize}
\item $|x|\in W^{1,p}(-1,1)$，弱导数为 $\operatorname{sgn}x$；Heaviside 的导数为 Dirac，故不属于这些 $W^{1,p}$。
\item $\Omega=(0,1)$ 上的常数一（$1\le p<\infty$）属于 $C^\infty(\Omega)\cap W^{1,p}$，但不属于 $W_0^{1,p}$；下一单元用零边界 Poincar\'e 不等式完成判别。
\end{itemize}

\planitem{输出与练习}
输出 Sobolev 完备性、AC/ACL 代表、乘积法则、卷积求导和 Meyers--Serrin。练习要求证明弱梯度为零的函数在每个连通分量上 a.e. 为常数、核算截断函数的弱导数、为带孔开集写出一组合法的局部卷积尺度，并由测试函数定义重证 $W^{1,1}$ 的双边卷积求导公式。

\unithead{4.3.4}{Poincar\'e、Friedrichs 与 Lipschitz 延拓}

\planitem{驱动问题}
弱导数控制函数本身时，常数必须怎样归一化？把域内函数送到全空间时，边界几何怎样进入算子范数？

\planitem{Poincar\'e 与 Friedrichs}
先在有界凸开集上证明
\[
\|u-u_\Omega\|_{L^p(\Omega)}
\le C(n,p,\Omega)\|\nabla u\|_{L^p(\Omega)},
\quad 1\le p<\infty,
\]
其中 $u_\Omega=|\Omega|^{-1}\int_\Omega u$。证明由 4.3.3 的线段 AC 公式、两次积分和核 $|x-y|^{1-n}$ 的可积性完成。对有界开集的 $W_0^{1,p}$ 证明
\[
\|u\|_p\le C_0(n,p,\Omega)\|\nabla u\|_p,
\]
先对 $C_c^\infty$ 沿一条坐标方向积分，再取闭包。由此得到梯度范数在 $W_0^{1,p}$ 上与完整范数等价。

\planitem{延拓定理的依赖闭合}
先定义有界 Lipschitz 域：对每个 $z\in\partial\Omega$，存在刚性坐标、半径 $r_z>0$ 和
Lipschitz 函数 $\phi_z$，使得在该坐标柱 $B'(0,r_z)\times(-r_z,r_z)$ 内
$\Omega=\{(x',x_n):x_n>\phi_z(x')\}$；用有限个这样的图覆盖边界。
对 $1\le p<\infty$ 证明图坐标变换引理：若
$\Phi(x',s)=(x',s+\phi(x'))$ 且 $\phi$ Lipschitz，则 $u\mapsto u\circ\Phi$ 在 $W^{1,p}$ 上有界，并给出 a.e. 链式法则。证明先由 $\phi$ 在几乎每条坐标线上的一维 AC 性识别其弱偏导，再依次处理光滑 $u$、Fubini 换元与 Meyers--Serrin 极限。

在半空间图上作反射延拓，乘单位分解后拼合，对 $1\le p<\infty$ 得到有界线性算子
\[
E:W^{1,p}(\Omega)\longrightarrow W^{1,p}(\R^n),
\qquad Eu=u\ \text{a.e. 于 }\Omega,
\]
并以一个在 $\overline\Omega$ 邻域恒为一的截断使 $Eu$ 支撑于固定球。零延拓作为 $W_0^{1,p}$ 的专用构造单独证明。

\planitem{主例}
在长盒 $\Omega_L=(0,L)\times(0,1)^{n-1}$ 上取 $u=x_1-L/2$，完整计算
\[
\|u\|_p^p=\frac{L^{p+1}}{2^p(p+1)},
\qquad \|\nabla u\|_p^p=L,
\]
故 Poincar\'e 常数至少按 $L$ 线性增长。再在半空间中取
$u(x',x_n)=\eta(x')\chi(x_n)$，其中 $\eta\in C_c^\infty(\R^{n-1})$、$\chi\in C_c^\infty(\R)$ 且 $\chi(0)=1$。对其朴素零延拓 $\widetilde u$ 直接算出
\[
\langle\partial_n\widetilde u,\varphi\rangle
=\int_{\{x_n>0\}}(\partial_n u)\varphi\dd x
+\int_{\R^{n-1}}\eta(x')\varphi(x',0)\dd x',
\]
由第二项直接看见边界上的额外分布，再与反射延拓及 $W_0^{1,p}$ 的零延拓并排比较。

\planitem{输出与练习}
输出域假设显式的 Poincar\'e/Friedrichs 不等式、图坐标引理、反射延拓与零延拓判据。练习先约定
$B_0\Subset\Omega$ 且“关于球 $B_0$ 星形”意为
$[x,y]\subset\Omega$ 对所有 $x\in\Omega$、$y\in B_0$ 成立，再比较凸域与此类星形域的常数；另对一个折线 Lipschitz 图手算反射后的弱梯度，并用常数函数和 Poincar\'e 不等式判定其不属于 $W_0^{1,p}(0,1)$。

\unithead{4.3.5}{Sobolev--Morrey 嵌入与端点诊断}

\planitem{驱动问题}
延拓把问题送到全空间后，缩放关系允许哪些目标指数？临界指数处是连续嵌入、紧嵌入还是端点失败？

\planitem{H\"older 记号}
对紧集 $K$ 与 $0<\alpha\le1$ 定义
\[
[u]_{C^{0,\alpha}(K)}=\sup_{x\ne y\in K}
\frac{|u(x)-u(y)|}{|x-y|^\alpha},
\qquad
\|u\|_{C^{0,\alpha}(K)}=\|u\|_{\infty,K}+[u]_{C^{0,\alpha}(K)},
\]
并记 $C^{0,\alpha}(K)$ 为该范数有限的连续函数空间；这里 $C(K)$ 表示紧集上的连续函数，约定 $C^{0,0}(K)=C(K)$，只在
$\alpha>0$ 时使用 Hölder 半范数。后文的 Morrey 估计、Arzel\`a--Ascoli 和
紧嵌入均以这一记号为准。

\planitem{嵌入定理}
设 $\Omega$ 为有界 Lipschitz 域、$1\le p<\infty$：
\begin{enumerate}
\item $p<n$ 时，$p^*=np/(n-p)$，证明
$W^{1,p}(\Omega)\hookrightarrow L^q(\Omega)$，$1\le q\le p^*$；核心估计先在 $C_c^\infty(\R^n)$ 证明，再由延拓与稠密性推广。
\item $p=n\ge2$ 时，证明嵌入每个有限 $L^q$，并给出常数对 $q$ 的依赖；随后由显式临界序列检验 $L^\infty$ 端点。
\item $p>n$ 时，证明 Morrey 估计和唯一 $C^{0,1-n/p}(\overline\Omega)$ 代表。
\item $n=1,p=1$ 单独由 AC 代表得到 $W^{1,1}(a,b)\hookrightarrow L^\infty(a,b)$。
\end{enumerate}

\planitem{缩放与临界例}
若 $1\le p<n$ 且 $0\ne\eta\in C_c^\infty(B(0,1))$，取
\[
u_\lambda(x)=\lambda^{(n-p)/p}\eta(\lambda x).
\]
逐项计算 $\|\nabla u_\lambda\|_p$ 与 $\|u_\lambda\|_{p^*}$ 不随 $\lambda$ 变化，而 $L^q$ 范数的幂次在 $q=p^*$ 处恰好换号。对每个固定 $q>p^*$，选择
$n/q\le a<n/p-1$；再固定
$\chi\in C_c^\infty(B(0,1))$、$\chi\equiv1$ 于 $B(0,1/2)$，令
$v_a(0)=0$、$v_a(x)=\chi(x)|x|^{-a}$（$x\ne0$），逐项核对 $v_a\in W^{1,p}$ 而
$v_a\notin L^q$。

当 $p=n\ge2$、$k\ge2$ 时，取径向函数
\[
u_k(r)=
\begin{cases}
(\log k)^{(n-1)/n},&0\le r\le k^{-1},\\
(\log k)^{-1/n}\log(1/r),&k^{-1}<r<1,\\
0,&r\ge1.
\end{cases}
\]
核对连接处连续，并用极坐标算出
\[
\|\nabla u_k\|_n^n
=|\mathbb S^{n-1}|(\log k)^{-1}
\int_{1/k}^1\frac{\dd r}{r}
=|\mathbb S^{n-1}|,
\qquad
\|u_k\|_\infty=(\log k)^{(n-1)/n}\to\infty.
\]
再分别估计中心球与外环，得到 $\sup_k\|u_k\|_n<\infty$，从而给出临界 $W^{1,n}\not\hookrightarrow L^\infty$ 的完整诊断。

\planitem{输出与练习}
输出次临界、临界和 Morrey 三段嵌入及其端点例。练习要求用缩放判定候选嵌入的唯一可能指数、计算 $u_k$ 的 $L^q$ 范数、构造一维 $W^{1,1}$ 代表并直接核对其 $L^\infty$ 控制。

\unithead{4.3.6}{Kolmogorov--Riesz 与 Rellich：唯一的函数空间紧性高潮}

\planitem{驱动问题}
$L^p$ 有界只能限制总质量，不能阻止质量压缩、高频振荡或向无穷远移动。哪些可直接检验的条件恰好把无限维族压成有限网？

\planitem{Kolmogorov--Riesz 定理}
对 $1\le p<\infty$，证明 $\mathcal F\subset L^p(\R^n)$ 相对紧当且仅当：
\begin{enumerate}
\item $\sup_{f\in\mathcal F}\|f\|_p<\infty$；
\item $\lim_{h\to0}\sup_{f\in\mathcal F}\|\tau_hf-f\|_p=0$；
\item $\lim_{R\to\infty}\sup_{f\in\mathcal F}
\|f\mathbf1_{\{|x|>R\}}\|_p=0$。
\end{enumerate}
充分性证明固定大立方体并按小网格取分块平均；用 Jensen 把分块误差控制成小平移的平均；这些分块落在有限维空间的有界集内，故有有限网；再加尾部误差。必要性从相对紧闭包的有限 $\eps$-网出发，把单个函数的平移连续与截尾统一到全族。

\planitem{紧性三种障碍}
固定 $0\ne\eta\in C_c^\infty(B(0,1))$，逐项计算：
\begin{itemize}
\item 集中：$f_k=k^{n/p}\mathbf1_{B(0,1/k)}$ 的 $L^p$ 范数恒定、尾部受控；给定任意小 $h\ne0$，选 $k$ 使 $1/k<|h|/3$，两支撑分离，从而一致平移模失败。
\item 逃逸：$f_k=\eta(\cdot-ke_1)$ 的范数和平移模完全一致，尾部条件失败。
\item 高频振荡：$f_k(x)=\eta(x)e^{ikx_1}$ 的范数和尾部受控；取
$h_k=\pi e_1/k$，利用 $e^{ik h_{k,1}}=-1$ 与 $L^p$ 平移连续性算出
\[
\|\tau_{h_k}f_k-f_k\|_p\longrightarrow2\|\eta\|_p,
\]
故一致小平移模失败。
\end{itemize}
由此在章首和本节统一表述：平移模排除集中与振荡，尾部条件排除逃逸。

\planitem{Rellich--Kondrachov}
设 $\Omega$ 为有界 Lipschitz 域。证明：
\begin{enumerate}
\item $1\le p<n$ 时，$W^{1,p}(\Omega)\Subset L^q(\Omega)$ 对 $1\le q<p^*$；
\item $p=n$ 时，对每个有限 $q$ 紧嵌入 $L^q$；
\item $p>n$ 时，对每个 $0<\beta<1-n/p$ 紧嵌入 $C^{0,\beta}(\overline\Omega)$。
\end{enumerate}
先由延拓与线段公式证明
$\|\tau_hu-u\|_p\le |h|\|\nabla u\|_p$，固定支撑给出尾部条件，调用 Kolmogorov--Riesz 得 $L^p$ 强收敛子列；再用 Sobolev 界和 $L^p$ 插值升级到次临界 $L^q$。$p>n$ 情形先以 Morrey 等度连续和 Arzel\`a--Ascoli 得到一致收敛，再用共同 $C^{0,1-n/p}$ 界与 H\"older 半范数插值得到每个 $C^{0,\beta}$（$\beta<1-n/p$）中的收敛。

\planitem{尖锐性例}
若 $1\le p<n$、沿用上面的 $0\ne\eta\in C_c^\infty(B(0,1))$ 且 $x_0\in\Omega$，对充分大的 $k$ 令
\[
u_k(x)=k^{(n-p)/p}\eta(k(x-x_0)).
\]
计算 $\|\nabla u_k\|_p$ 与 $\|u_k\|_{p^*}$ 均恒定，而
$\|u_k\|_p=k^{-1}\|\eta\|_p$。若临界 $L^{p^*}$ 存在强收敛子列，其极限由依测度收敛只能为零，与恒定范数矛盾。

\planitem{输出与练习}
输出完整 KR 充要判据和 Rellich 紧嵌入。练习要求为三种障碍逐项核对另外两项条件、用分块平均实际构造给定误差的有限网、比较临界缩放在 $L^p$ 与 $L^{p^*}$ 中的行为。

\subsection*{4.4　演化链：从核的复合到生成元}

\unithead{4.4.1}{卷积代数、测度卷积与有限测度平滑}

\planitem{驱动问题}
卷积不只是一次光滑操作：重复卷积会组合平移、叠加测度，并产生随时间复合的算子。先要确定这些运算在哪个空间闭合，以及函数卷积与测度卷积怎样兼容。

\planitem{平移定理的统一登记}
将 4.3.1 的逐截面论证记为
\texttt{theorem:4-4-1-1}：对 $1\le p<\infty$，$\tau_hf\to f$ 于 $L^p(\R^n)$，
并且对有限测度空间 $S$，相应的平移在 $L^p(S\times\R)$ 中强连续。此处把两种
表述并列核对，并说明它们分别为卷积结合律和第六章时间变量平移提供输入。

\planitem{$L^1$ 卷积代数}
利用 4.3.1 已经固定的可积性，证明
\[
f*g=g*f,\qquad (f*g)*h=f*(g*h),\qquad
\|f*g\|_1\le\|f\|_1\|g\|_1.
\]
结合律的证明把三重绝对积分先写出，再应用 Fubini 和线性换元。对明确代表证明
\[
\ess\supp(f*g)
\subset\overline{\ess\supp f+\ess\supp g}.
\]
证明 $L^1(\R^n)$ 没有卷积单位元：若 $e\in L^1$ 为单位，取一列 $\eps_m\downarrow0$，则
$e*\rho_{\eps_m}=\rho_{\eps_m}$ 在 $L^1$ 等价类意义下成立。取这些等式成立集合的可数交，再与 $e$ 的 Lebesgue 点集相交；在所得共同满测集上固定 $x\ne0$ 并令 $m\to\infty$，左端趋于 $e(x)$，右端最终为零，故 $e=0$ a.e.，与单位性质矛盾。有限测度卷积代数的单位元 $\delta_0$ 随后成为扩张动机。

\planitem{测度卷积}
对有限 Borel 测度 $\mu,\nu$，记 $\|\mu\|_{TV}=|\mu|(\R^n)$，并以加法映射
$S(x,y)=x+y$ 定义
\[
\mu*\nu=S_*(\mu\otimes\nu).
\]
先对正测度，再经 Jordan/极分解扩张到有限有符号或复测度，证明
\[
\|\mu*\nu\|_{TV}\le\|\mu\|_{TV}\|\nu\|_{TV},
\quad \delta_0*\mu=\mu,
\]
以及交换、结合。若 $\mu=fm$、$\nu=gm$，用 Fubini 与换元证明
$(fm)*(gm)=(f*g)m$；若一项有密度，则写出函数--测度卷积
\[
(\mu*\phi)(x)=\int\phi(x-y)\dd\mu(y),\qquad
\|\mu*\phi\|_1\le\|\mu\|_{TV}\|\phi\|_1.
\]

\planitem{有限测度的平滑恢复}
若 $\mu$ 为有限测度，$\phi\in L^1(\R^n)$、$\phi\ge0$、$\int\phi=1$，并令
$\phi_\eps(x)=\eps^{-n}\phi(x/\eps)$。先约定反射核
$\check\phi(x)=\phi(-x)$，再对每个 $g\in C_b(\R^n)$ 证明
\[
\int g(x)(\mu*\phi_\eps)(x)\dd x
=\int(g*\check\phi_\eps)(y)\dd\mu(y)
\longrightarrow\int g\dd\mu.
\]
第一步由 Fubini，第二步对每个 $y$ 用 $g$ 在 $y$ 的连续性，再以
$\|g*\check\phi_\eps\|_\infty\le\|g\|_\infty$ 对 $|\mu|$ 作有界收敛。该命题补齐第五章用高斯平滑有限测度再撤去平滑的接口。

\planitem{Gaussian 例}
先定义有限测度的特征函数
$\widehat\mu(\xi)=\int e^{\,i\xi\cdot x}\dd\mu(x)$，再由密度完成平方证明非退化
Gaussian 的卷积公式；固定半正定矩阵 $Q$ 的对称平方根 $Q^{1/2}$，定义
$\gamma_Q=(Q^{1/2})_*\gamma_I$。由标准 Gaussian 密度的正交不变性（在此处直接证明）
可验证任意 $A$ 满足 $AA^{\mathsf T}=Q$ 时 $A_*\gamma_I$ 给出同一测度；
其特征函数计算为 $\exp(-\tfrac12\xi^{\mathsf T}Q\xi)$。乘积测度经
$(u,v)\mapsto Au+Bv$ 的推前给出
\[
\gamma_Q*\gamma_R=\gamma_{Q+R},
\]
包含退化协方差。章际接口由加法推前和这一恒等式构成；第五章把它解释为独立随机变量之和的分布。

\planitem{输出与练习}
输出函数/测度卷积代数、有限测度测试恢复和 Gaussian 闭包。练习要求由两原子测度手算卷积、验证三角形密度与加法推前一致、证明总质量与一阶矩在可积条件下的卷积公式。

\unithead{4.4.2}{近似恒等的范数恢复、逐点恢复与诊断}

\planitem{驱动问题}
哪些核族在尺度趋零时恢复恒等算子？范数恢复与每个 Lebesgue 点恢复需要哪些不同的控制，而每一项假设失效时会出现什么可计算的现象？

\planitem{近似恒等的定义与范数恢复}
称 $(\phi_\eps)\subset L^1(\R^n)$ 为近似恒等，若
\[
\int\phi_\eps=1,\qquad
\sup_\eps\|\phi_\eps\|_1<\infty,\qquad
\int_{|x|>\delta}|\phi_\eps(x)|\dd x\longrightarrow0
\quad(\delta>0).
\]
证明对 $1\le p<\infty$ 的 $f\in L^p$ 有
$\|\phi_\eps*f-f\|_p\to0$，对 $f\in C_0$ 有一致收敛。证明统一写成
\[
\phi_\eps*f-f=
\int\phi_\eps(y)(\tau_yf-f)\dd y
\]
并分近区、远区；$C_0$ 的一致平移连续由第一章的紧性论证给出。该结论沿用标签 \texttt{theorem:4-4-2-1}。

\planitem{Lebesgue 点恢复}
先约定：非负函数 $\Psi$ 称为径向递减，是指存在非增函数
$\psi:[0,\infty)\to[0,\infty)$ 使 $\Psi(x)=\psi(|x|)$。若固定形状 $\phi$
满足 $\int\phi=1$ 且 $|\phi|\le\Psi$，其中 $\Psi\in L^1$ 径向递减，则证明
\[
\sup_{\eps>0}|\phi_\eps*f(x)|
\le C_n\|\Psi\|_1Mf(x),
\qquad
\phi_\eps*f(x)\to f(x)
\]
在 $f\in L^p(\R^n)$（$1\le p\le\infty$）的每个 Lebesgue 点成立。证明用二进环层把核积分化成球平均的可和组合；内层由 Lebesgue 点的平均振荡控制，外层由 $Mf(x)$ 与 $\Psi$ 的可积尾部控制。径向大函数是一条可直接核验的充分条件。该结论沿用标签 \texttt{theorem:4-4-2-2}。

\planitem{算子范数的 Baire 诊断}
在反例计算前证明所需的一致有界特例：若 $X,Y$ 为 Banach 空间、
$\mathcal B(X,Y)$ 表示有界线性算子空间，$S_k\in\mathcal B(X,Y)$，且对每个 $x\in X$ 都有
$\sup_k\|S_kx\|<\infty$，则 $\sup_k\|S_k\|<\infty$。
令 $E_N=\{x:\sup_k\|S_kx\|\le N\}$；这些闭集覆盖 $X$，Baire 定理给出某个
$E_N$ 含有一个球，取球内两点之差并按半径缩放即得算子范数的一致界。
因此若 $\sup_k\|S_k\|=\infty$，就存在 $x\in X$ 使 $\sup_k\|S_kx\|=\infty$。
对卷积算子 $S_kf=\phi_k*f$，先由 Young 上界和 4.3.1 的单位质量 mollifier
$\rho_\delta$ 的测试（$\|\phi_k*\rho_\delta\|_1\to\|\phi_k\|_1$）得到
$\|S_k\|_{L^1\to L^1}=\|\phi_k\|_1$，再将这一推论用于下一项诊断。

\planitem{紧跟假设的四个诊断例}
\begin{enumerate}
\item 固定标准核 $\rho$，令 $\phi_\eps=a\rho_\eps$（$a\ne1$），并取
$f(x)=e^{-|x|^2}$。直接算出 $\int\phi_\eps=a$，且
$(\phi_\eps*f)(0)\to a\ne f(0)=1$；一般极限相应为 $af$。
\item 取 $\rho_\eps(\cdot-a)$、$a\ne0$，其质量和 $L^1$ 范数保持，但质量集中在 $a$，卷积极限为相应平移；这精确对应尾部集中条件。
\item 取 $\rho,\psi\in C_c^\infty(B(0,1))$，满足 $\rho\ge0$、$\int\rho=1$、$\int\psi=0$、$\psi\ne0$，令
$\rho_k(x)=k^n\rho(kx)$、$\psi_k(x)=k^n\psi(kx)$ 及
$\phi_k=\rho_k+k\psi_k$。计算
$\int\phi_k=1$ 且 $\|\phi_k\|_1\to\infty$；卷积算子在 $L^1$ 上的范数由近似点质量给出同阶下界，再由上一项的算子 Baire 结论得到一个 $f\in L^1$，使 $\phi_k*f$ 不能在 $L^1$ 中收敛。
\item 令 $r_k=2^{-k}$、$\ell_k=2^{-2k}$、
$E=\bigcup_k(r_k,r_k+\ell_k)$，则 $0$ 是 $\mathbf1_E$ 的密度零点；而
\[
\phi_k(y)=\ell_k^{-1}\mathbf1_{(-r_k-\ell_k,-r_k)}(y)
\]
满足近似恒等的三项范数条件，却有
$(\phi_k*\mathbf1_E)(0)=1$。这说明范数恢复条件与逐 Lebesgue 点恢复条件的层级差异，并定位径向大函数条件在逐点结论中的作用。
\end{enumerate}

最后回访 E1：把 $f=\mathbf1_{B(0,1)}$ 代入径向大函数估计，分别在 $|x|\le2$ 与 $|x|>2$ 使用 4.1.2 已算出的 $Mf(x)$，得到核最大函数的有界区控制和 $|x|^{-n}$ 远场衰减。由此把尺度链的极大估计落实为本单元的逐点恢复工具。

\planitem{输出与练习}
输出近似恒等的范数恢复、径向大函数条件下的 Lebesgue 点恢复以及四项独立诊断。练习要求逐项验证一个变号近似恒等、比较 E2 的对称区间平均与偏心核取样、构造质量一且尾部集中但 $L^1$ 范数增长的另一核列。

\unithead{4.4.3}{热核、Poisson 核与正时间方程}

\planitem{驱动问题}
哪些近似恒等还具有参数相加的复合律？核本身满足的微分方程怎样传给粗糙初值，并在正时间产生经典解？

\planitem{热核}
定义并沿用标签 \texttt{def:4-4-2-2}
\[
p_t(x)=(4\pi t)^{-n/2}e^{-|x|^2/(4t)},\qquad t>0.
\]
通过 Gaussian 完成平方逐项证明
\[
\int p_t=1,\qquad p_s*p_t=p_{s+t},\qquad
\partial_tp_t=\Delta p_t,
\]
以及 $\int x_jp_t=0$、$\int x_j^2p_t=2t$，并核对
$\int_{|x|>\delta}p_t(x)\,\dd x\to0$（$\delta>0$）。因此
$H_tf=p_t*f$ 保正、在各 $L^p$ 上收缩，并在 $1\le p<\infty$ 与 $C_0$ 上构成强连续半群。第六章采用
$q_t=p_{t/2}$，其生成元系数相应为 $\tfrac12\Delta$。

\planitem{Poisson 核}
定义（$\Gamma$ 为 Euler Gamma 函数）
\[
P_t(x)=c_n\frac{t}{(t^2+|x|^2)^{(n+1)/2}},
\qquad c_n=\frac{\Gamma((n+1)/2)}{\pi^{(n+1)/2}}.
\]
用 Beta--Gamma 积分核对质量一和径向递减，直接微分证明
$(\Delta_x+\partial_t^2)P_t=0$。半群律由自足的次级化公式
\[
P_t(x)=\frac{t}{2\sqrt\pi}
\int_0^\infty e^{-t^2/(4s)}s^{-3/2}p_s(x)\dd s
\]
推出。具体令
\[
\eta_t(s)=\frac{t}{2\sqrt\pi}e^{-t^2/(4s)}s^{-3/2}
\mathbf1_{(0,\infty)}(s),
\]
先用 Gaussian 积分证明
\[
\int_0^r\eta_t(s)\eta_u(r-s)\dd s=\eta_{t+u}(r),
\]
再以 Tonelli 和热核半群律得到 $P_t*P_u=P_{t+u}$。在 $\R$ 上完整计算
\[
(P_t*\mathbf1_{[-1,1]})(x)=\frac1\pi
\left(\arctan\frac{1-x}{t}+\arctan\frac{1+x}{t}\right),
\]
读出内部、外部和边界三种极限。
由 $P_t\ge0$、$\int P_t=1$ 以及缩放给出的
$\int_{|x|>\delta}P_t(x)\,\dd x\to0$，按近似恒等定义补出
$P_t*f$ 在 $1\le p<\infty$ 的 $L^p$ 范数收缩和 $C_0$ 上的一致强连续性；
$L^\infty$ 端点的结论取为 Lebesgue 点逐点恢复；范数恢复的范围为 $1\le p<\infty$ 与 $C_0$。

\planitem{核作用后的方程}
对 $1\le p\le\infty$ 与 $f\in L^p(\R^n)$，证明
\[
u(x,t)=H_tf(x),\qquad v(x,t)=P_t*f(x)
\]
在每个正时间带 $\R^n\times[a,b]$（$0<a<b<\infty$）上关于 $(x,t)$ 光滑，并分别满足
\[
\partial_tu=\Delta_xu,
\qquad
(\Delta_x+\partial_t^2)v=0.
\]
证明先逐项核对热核与 Poisson 核的所有时空导数在固定正时间带具有可积大函数，再用 Young/H\"older、DCT 和卷积求导交换微分。对 $1\le p<\infty$ 及 $C_0$ 初值，4.4.2 给出 $t\downarrow0$ 的强恢复；Lebesgue 点处的恢复由径向递减大函数条件给出。这一结论为 4.4.5 的一般 $L^2$ 热能量恒等式提供正时间正则性。

\planitem{输出与练习}
输出热/Poisson 半群、核方程、正时间经典解与初值恢复。练习要求从热核二阶矩推导小时间 Taylor 余项、逐个验证核导数的正时间带支配、计算 Poisson 核在区间端点的 $1/2$ 极限。

\unithead{4.4.4}{矩阵实验、$C_0$-半群与生成元}

\planitem{驱动问题}
热核的时间相加、平移的参数相加和矩阵指数的指数律具有同一结构。抽去公式后，哪些性质足以保证轨道可连续观察，并产生一个可能无界的无穷小算子？

\planitem{有限维实验}
对
\[
A=\begin{pmatrix}-1&1\\1&-1\end{pmatrix}
\]
直接算出
\[
e^{tA}=\frac12
\begin{pmatrix}1+e^{-2t}&1-e^{-2t}\\1-e^{-2t}&1+e^{-2t}\end{pmatrix}.
\]
逐项验证正性、坐标和保持、半群律及 $\frac\dd{\dd t}e^{tA}=Ae^{tA}$。再对
$\operatorname{Re}\lambda>0$ 直接积分
$\int_0^\infty e^{-\lambda t}e^{tA}\dd t=(\lambda I-A)^{-1}$，作为 4.4.5 的有限维预演。

\planitem{定义与局部一致有界}
记 $\mathcal B(X)$ 为 $X$ 上有界线性算子的 Banach 空间，范数为
$\|T\|=\sup_{\|x\|\le1}\|Tx\|$。在 Banach 空间 $X$ 上定义
$C_0$-半群 $\{T_t\}_{t\ge0}\subset\mathcal B(X)$：
\[
T_0=I,\qquad T_{s+t}=T_sT_t,\qquad
\lim_{t\downarrow0}\|T_tx-x\|=0\quad(x\in X).
\]
若 $\|T_t\|\le1$ 对所有 $t\ge0$，称其为收缩半群；在带逐点序且含常数一的函数
Banach 空间中，保正并满足 $T_t\mathbf1=\mathbf1$ 称为 Markov 半群。无穷测度的
$L^p$（$p<\infty$）不含常数一时，相应的次 Markov 条件写成
$0\le f\le1\Rightarrow0\le T_tf\le1$ a.e.（在 $L^p\cap L^\infty$ 上理解）。
先由零时刻强连续得到：对每个固定 $x$，存在 $N$ 使
$\sup_{0\le s\le1/N}\|T_sx\|\le N$。令
$E_N=\{x:\sup_{0\le s\le1/N}\|T_sx\|\le N\}$，这些集合为闭集且 $X=\bigcup_NE_N$；第一章的 Baire 定理给出某个
$E_N$ 有内点，取其内点差集并缩放即可得到
$C_\delta=\sup_{0\le s\le\delta}\|T_s\|<\infty$（$\delta=1/N$）。再把
$t$ 分成有限多个不超过 $\delta$ 的时间段，推出
$\sup_{0\le t\le T}\|T_t\|<\infty$，并由半群律推出存在
$M\ge1$、$\omega\in\R$ 使
\[
\|T_t\|\le Me^{\omega t}.
\]
同时利用这一局部算子界与半群律证明全部时间的轨道连续。用对角半群
\[
(T_ta)_k=e^{-tk^2}a_k\quad(a\in\ell^2(\mathbb Z))
\]
计算 $\|T_t-I\|=1$（$t>0$），明确区分逐向量强连续与算子范数连续。

\planitem{生成元}
定义并沿用标签 \texttt{def:4-4-3-2}
\[
D(A)=\left\{x\in X:\lim_{t\downarrow0}\frac{T_tx-x}{t}
\text{ 在 }X\text{ 中存在}\right\},
\qquad Ax=\lim_{t\downarrow0}\frac{T_tx-x}{t}.
\]
对
\[
x_h=\frac1h\int_0^hT_sx\dd s
\]
证明 $x_h\in D(A)$、$x_h\to x$ 且
$Ax_h=(T_hx-x)/h$，从而 $D(A)$ 稠密。随后对 $x\in D(A)$ 建立
\[
T_tx\in D(A),\qquad AT_tx=T_tAx,\qquad
\frac\dd{\dd t}T_tx=T_tAx,
\]
以及轨道积分公式
\[
T_tx-x=\int_0^tT_sAx\dd s.
\]
最后令 $x_m\to x$、$Ax_m\to y$，在轨道积分公式中取极限，再令 $t\downarrow0$，得到 $x\in D(A)$ 且 $Ax=y$，从而证明生成元闭。
轨道交换结论沿用标签 \texttt{proposition:4-4-3-1}。

\planitem{无界生成元例}
对上述对角半群逐坐标识别
\[
(Aa)_k=-k^2a_k,\qquad
D(A)=\left\{a\in\ell^2(\mathbb Z):
\sum_kk^4|a_k|^2<\infty\right\}.
\]
先证右侧包含于定义域，再用 Fatou 从差商的 $\ell^2$ 收敛推出反向包含；由标准基向量看见 $A$ 无界。

\planitem{输出与练习}
输出增长界、稠密闭生成元、轨道微分和对角模型。练习要求比较两个不同矩阵指数的生成元、验证右移半群在 $L^\infty$ 上强连续失败、由 $x_h$ 公式估计图范数。

\unithead{4.4.5}{具体生成元、能量与 Laplace--Bochner 预解式}

\planitem{驱动问题}
生成元只由零时刻差商定义，怎样从整个时间轨道反推出 $(\lambda I-A)^{-1}$？在平移和热核两个模型中，这个无穷小算子具体是什么？

\planitem{平移与热生成元}
对 $1\le p<\infty$，令 $T_tf(x)=f(x+t)$ 作用于 $L^p(\R)$。证明
\[
D(A)=W^{1,p}(\R),\qquad Af=f'.
\]
$W^{1,p}$ 到定义域的方向用 AC 代表和积分型 Minkowski；反向把强差商极限放入测试函数配对，识别为弱导数。若采用 $f(x-t)$ 的平移约定，正文同步把生成元写成 $-f'$。

分别在 $X=L^p(\R^n)$（$1\le p<\infty$）与 $X=C_0(\R^n)$ 上取热半群；4.4.3 已证明它们是 $C_0$-半群。对每个 $f\in C_c^\infty(\R^n)$ 证明
$f\in D(A)$ 且 $Af=\Delta f$：把 $p_t*f-f$ 写成二阶 Taylor 公式的 Gaussian 平均，一阶矩消失，二阶矩给出系数。由此建立稠密测试类上的生成元识别；随后的能量恒等式固定在 $X=L^2(\R^n)$ 上。
这里 $C_c^\infty(\R^n)$ 在 $C_0(\R^n)$ 中的稠密性由第一章的 $C_c$ 截断和
4.3.1 的一致卷积逼近逐步得到，使两种状态空间的“稠密测试类”含义一致。

\planitem{热能量}
先对光滑数据 $u(t)=H_tf$ 写出复 $L^2$ 中的公式
\[
\frac\dd{\dd t}\|u(t)\|_2^2
=2\operatorname{Re}\langle\Delta u(t),u(t)\rangle
=-2\|\nabla u(t)\|_2^2.
\]
对一般 $f\in L^2$，先在任意 $0<a<b$ 上把同一恒等式应用于已经光滑化的 $H_af$，以截断分部积分消去无穷远边界；再用 $H_af\to f$ 的强连续性并令 $a\downarrow0$，得到
\[
\|H_bf\|_2^2+2\int_0^b\|\nabla H_sf\|_2^2\dd s
=\|f\|_2^2.
\]

\planitem{预解式}
先在复 Banach 空间上定义闭算子 $A$ 的预解集与 $R(\lambda,A)=(\lambda I-A)^{-1}$。若
$\|T_t\|\le Me^{\omega t}$ 且 $\operatorname{Re}\lambda>\omega$，定义
\[
R_\lambda x=\int_0^\infty e^{-\lambda t}T_tx\dd t.
\]
由 Bochner 积分证明
\[
\|R_\lambda\|\le\frac{M}{\operatorname{Re}\lambda-\omega}.
\]
变量换元计算
$[T_hR_\lambda x-R_\lambda x]/h$ 的极限，得到
\[
R_\lambda X\subset D(A),\qquad
(\lambda I-A)R_\lambda=I.
\]
对 $x\in D(A)$ 积分
\[
\frac\dd{\dd t}\bigl(e^{-\lambda t}T_tx\bigr)
=-e^{-\lambda t}T_t(\lambda x-Ax)
\]
并令上端趋于无穷，得到
$R_\lambda(\lambda I-A)x=x$。最后由两个逆恒等式推出
\[
R_\lambda-R_\mu=(\mu-\lambda)R_\lambda R_\mu.
\]
这一完整结论沿用标签 \texttt{proposition:4-4-3-4}。

\planitem{Yosida 近似}
由预解式定义
\[
A_\lambda=\lambda AR_\lambda
=\lambda^2R_\lambda-\lambda I
\]
（实参数 $\lambda>\max\{\omega,0\}$），证明
$\lambda R_\lambda x\to x$ 对每个 $x\in X$ 成立，并对
$x\in D(A)$ 证明 $A_\lambda x\to Ax$。标量半群
$T_t=e^{-ct}I$ 与对角半群分别给出可逐坐标核对的公式。

\planitem{输出与练习}
输出两个具体生成元、热能量、预解式的两侧逆、范数界和预解恒等式。练习要求对矩阵实验直接验证全部公式、计算平移半群预解式为单侧指数卷积、比较热核 $p_t$ 与第六章归一化 $p_{t/2}$ 的生成元系数。

\section{原第四章覆盖与内容归属}

下表按原第四章正文的顺序列出全部核心内容。合并发生在重复证明和同一例子的重复计算处；正式结果、关键反例、图示承担的数学关系以及练习能力目标均有明确落点。

\begin{longtable}{p{0.17\textwidth}p{0.36\textwidth}p{0.37\textwidth}}
\toprule
原章单元 & v2 落点 & 修订后的展开方式\\
\midrule
章首导言 & 修订目标、四条链、例题与诊断账本 & 以“有限观测--统一估计--极限--识别”为共同机制，并逐章列出真实输入\\
4.1.1 极大函数 & 4.1.1--4.1.2 & 局部平均的合法定义与可测性先行；$5r$ 选择、弱 $(1,1)$ 和强 $(p,p)$ 分层证明；重叠例计算集中于 4.1.1，4.1.2 调用其结论\\
4.1.2 Lebesgue 点 & 4.1.3 & 连续类、$L^1$ 逼近、极大坏集三步展开；密度点、精确代表、Smith--Volterra 对照完整保留；Riemann 判据调用第三章的已证版本\\
4.1.3 测度密度 & 4.1.4 & 函数密度、奇异测度、一般有符号测度及相对密度依次推进；Radon 正则性先局部有限化再使用\\
4.2.1 BV & 4.2.1--4.2.2 & 总变差、Jordan、BV 的代数与乘积估计、规范代表、BV--测度对应与单调/BV a.e. 可微全部保留；左端原子由二元组统一\\
4.2.2 Stieltjes & 4.2.2 & 三类可算积分、网和、算子范数、共同跳点和分部积分形成一条闭合证明链；向量积分直接调用第三章 Bochner 接口\\
4.2.3 AC 与 FTC & 4.2.3--4.2.4 & 从第三章 AC 积分表示推进到密度识别、变差等式、$C^1$ 换元、曲线，再由面积公式完成一般 Lipschitz 链式法则、三分解、indicatrix 和 Banach--Zarecki\\
4.3.1 mollifier & 4.3.1 & 卷积估计在首次使用前固定；全空间与局部恢复、任意开集上的 $C_c^\infty$ 稠密、边界层逐层展开\\
4.3.2 分布 & 4.3.2 & 测试空间连续性、测度嵌入、任意阶分布导数及混合偏导交换、光滑乘法、局部耗尽与光滑化、弱导数、闭性、$p>1$ 弱抽取及一维基本解例形成完整单元\\
4.3.3 Sobolev & 4.3.3--4.3.6 & 完备性、AC/ACL、Meyers--Serrin、卷积求导、Poincar\'e、延拓、嵌入、KR、Rellich 按依赖拆成四单元；Rellich 成为本章唯一的函数空间紧性高潮\\
4.4.1 卷积结构 & 4.3.1、4.4.1 & Young 首次用于光滑化时证明；交换/结合、支撑、无单位、测度卷积与 Gaussian 平滑在结构节汇总\\
4.4.2 核与恢复（含 Fourier/Brownian 解释） & 4.4.2--4.4.3；第六、七章 & 一般近似恒等、热核、Poisson 核、半群律与正时间方程在本章完整展开；Brownian 解释进入第六章的过程/转移核语境；Dirichlet/Fej\'er 核进入第七章，与 Fourier 部分和及核范数增长共同处理\\
4.4.3 半群 & 4.4.4--4.4.5 & 强连续、局部一致界、闭生成元、具体识别、能量、Yosida 与 Laplace--Bochner 预解式逐步证明\\
\bottomrule
\end{longtable}

\subsection{专题边界与后章衔接}

\begin{longtable}{p{0.25\textwidth}p{0.62\textwidth}}
\toprule
主题 & 在全卷中的自然位置\\
\midrule
Helly、Helly--Bray 与分布函数判据 & 4.2 完成端点二元组、Stieltjes 测度和连续测试积分；第七章在整轴质量控制的语境中完成选择定理与分布函数收敛\\
Dirichlet/Fej\'er 核与 Baire 发散 & 4.4 建立一般近似恒等、热核和 Poisson 核；第七章结合 Fourier 部分和、核范数增长及一致有界原理完成正反两面\\
有限状态 $Q$-矩阵与 Chapman--Kolmogorov & 4.4.4 的二维矩阵承担半群、正性和预解式的可算原型；第六章在转移概率与过程语境中建立完整有限状态理论\\
Laplace 基本解与 Newton 位势 & 4.3.2 以 $(|x|/2)''=\delta_0$ 完成分布基本解的可算接口；多维 Poisson 方程的存在、唯一性和正则性组成独立 PDE 发展线\\
一般微分基与强矩形基 & 4.1 证明欧氏球、立方体及有界离心率集合的完整结论；更一般的基由覆盖性质分类，适合与强极大函数共同发展\\
卷积支撑的精确端点定理 & 4.4.1 以 $\ess\supp(f*g)\subset\overline{\ess\supp f+\ess\supp g}$ 完成当前依赖链；精确支撑等式与复分析工具组成另一专题\\
\bottomrule
\end{longtable}

这些归属使第四章的新增内容都满足至少一个闭环：进入本章核心定理、成为第五至七章已有证明的直接输入，或完成原第四章已经展开的主干。

\section{依赖顺序与章际接口}

\subsection{十九个单元的依赖闭包}

\begin{longtable}{p{0.12\textwidth}p{0.37\textwidth}p{0.40\textwidth}}
\toprule
单元 & 进入本单元前已经具备 & 本单元新增并供后文调用\\
\midrule
4.1.1 & $L^1_{\rm loc}$、参数积分可测性、紧集有限覆盖 & 合法球平均、振荡可测性、$5r$ 选择\\
4.1.2 & 4.1.1、层蛋糕、Chebyshev & 极大弱型、弱 $L^1$、强 $L^p$\\
4.1.3 & 4.1.2、$C_c$ 稠密、积分绝对连续性 & Lebesgue 点、精确代表、偏心集合推论\\
4.1.4 & 4.1.1--4.1.3、RN/Lebesgue 分解、Radon 正则性 & 测度微分、单侧区间与相对密度\\
4.2.1 & 测度 Jordan 分解、Stieltjes 构造 & BV 代数与乘积估计、规范代表、端点二元组、变差测度\\
4.2.2 & 4.1.4、4.2.1、连续分离函数、Bochner 积分 & 单调/BV 可微、标量/Banach 值 Stieltjes 积分与范数\\
4.2.3 & 4.1.3、4.2.2、第三章 AC 表示 & FTC 闭合、变差等式、换元与曲线\\
4.2.4 & 4.2.1--4.2.3、测度三分解 & BV 三分解、indicatrix、AC 面积公式、Lipschitz 链式法则与 Banach--Zarecki\\
4.3.1 & Fubini/Minkowski、平移连续、$L^1$ 卷积 & Young、局部 mollifier、任意开集的 $C_c^\infty$ 稠密\\
4.3.2 & 4.2.1--4.2.2、4.3.1、测试函数微分 & 内部耗尽、测度分布嵌入、任意阶分布导数、光滑乘法、分布光滑化、弱导数的唯一与闭性\\
4.3.3 & 4.2.3、4.3.1--4.3.2、$L^p$ 完备 & $W^{1,p}$、ACL、卷积求导、Meyers--Serrin\\
4.3.4 & 4.3.3、光滑单位分解、凸域线段公式 & Poincar\'e、Friedrichs、Lipschitz 延拓\\
4.3.5 & 4.3.4、H\"older、径向积分与缩放 & 全空间/域上 Sobolev--Morrey 嵌入与端点诊断\\
4.3.6 & 4.3.4--4.3.5、有限维全有界、$L^p$ 完备 & Kolmogorov--Riesz、Rellich\\
4.4.1 & 4.3.1--4.3.2、乘积/推前测度、Fubini & 卷积代数、测度卷积、有限测度平滑\\
4.4.2 & 4.1.2--4.1.3、4.3.1、4.4.1 & 一般近似恒等、Lebesgue 点恢复与失败诊断\\
4.4.3 & 4.3.1、4.4.1--4.4.2 & 热/Poisson 卷积半群、核方程与正时间经典解\\
4.4.4 & 4.4.3、第一章 Baire、第三章 Bochner 积分 & $C_0$-半群、闭生成元、轨道微分\\
4.4.5 & 4.3.2--4.3.5、4.4.3--4.4.4 & 具体生成元、热能量、预解式、Yosida 近似\\
\bottomrule
\end{longtable}

依赖引用在正文中采用“定理名 + 实际 LaTeX 标签”。例如 Stieltjes 网和沿用第一章
\texttt{ex:1-1-1-2} 的带标记分割有向预序；半群一致有界引用
\texttt{theorem:1-4-1-1} 与 \texttt{theorem:1-4-1-3}；$C_0$ 完备性引用
\texttt{proposition:1-3-1-c0-complete}；距离函数分离与连续有限单位分解分别调用第一章的对应命题。4.3.3 所需的\emph{光滑}局部有限单位分解由该单元利用
\texttt{ex:1-3-1-smooth-bump} 另行构造。

\subsection{有向参数与选择账本}

第一次写 $r\downarrow0$ 时令 $D=(0,r_0)$，规定 $r\preceq s$ 当且仅当 $r\ge s$；正有理半径构成共尾子集。Stieltjes 分割以共同加细为有向关系；有限测试函数集以包含关系有向；半群预解式以有限积分上端 $R\uparrow\infty$ 有向。这样“同时满足有限批观测”在四条链中使用同一量词结构。

沿用第一章的选择标签：
\begin{itemize}
\item 紧集上的有限坏球覆盖、有限网格与有限分割只需有限选择，标记为 ZF；
\item 把第二可数开覆盖化为可数子族时标记 CC；按尺度层依赖既有选择作贪心抽取时标记 DC；
\item 相对测度微分的 Besicovitch 证明采用显式有限重叠算法，并逐层记录选球顺序与重叠常数；
\item 对可数稠密测试族逐层抽子列时标记 DC，并写明测试族的预先编号。
\end{itemize}
这些标签集中放在证明末的“基础账本”行，使正文计算保持连续。

\subsection{第五至七章的真实接口}

\begin{longtable}{p{0.27\textwidth}p{0.27\textwidth}p{0.35\textwidth}}
\toprule
第四章输出 & 后续位置 & 接口内容\\
\midrule
测度卷积与密度相容性 & 第五章独立和 & 加法推前先给测度恒等式，有密度时化为 $f*g$\\
有限测度的 Gaussian 平滑恢复 & 第五章特征函数唯一性 & 对 $C_b$ 测试函数撤去 Gaussian 平滑\\
乘积空间时间平移 & 第六章过程的时间平移 & 由 \texttt{theorem:4-4-1-1} 的有限测度乘积空间推论，在 $L^2(\Omega\times\R)$ 中获得强连续逼近\\
热核 $p_t$ 与 $q_t=p_{t/2}$ & 第六章 Brownian 转移核 & 半群律相同，生成元由 $\Delta$ 变为 $\tfrac12\Delta$\\
$C_0$ 与 $L^p$ 近似恒等 & 第六章 Feynman--Kac；第七章热/Fourier & 小时间强收敛与边界数据恢复\\
Lebesgue 点核恢复 & 第七章 Fourier 反演 & Gaussian 或其他受控核在典型点撤去正则化\\
$C_c^\infty$ 的 $L^p$ 稠密与卷积求导 & 第七章 Plancherel、$W^{1,1}$ 公式 & 先在好类证明，再以范数估计推广\\
闭生成元与 Laplace--Bochner 预解式 & 第七章半群与闭算子 & 两侧逆、增长界、预解恒等式均可直接引用\\
\bottomrule
\end{longtable}

\subsection{下游标签兼容表}

正文改写时，下列既有标签绑定到等价或更强的正式结论；构建全卷后逐一检查交叉引用。

\begin{longtable}{p{0.34\textwidth}p{0.50\textwidth}}
\toprule
既有标签 & v2 中绑定的结论\\
\midrule
\texttt{proposition:4-2-1-3} & $BV[a,b]$ 的向量空间与乘积封闭、齐次性、次可加性和乘积变差估计\\
\texttt{theorem:4-4-1-1} & $1\le p<\infty$ 的 $L^p$ 平移强连续及其有限测度乘积空间推论\\
\texttt{theorem:4-3-1-3} & 对任意开集 $\Omega\subset\R^n$，$C_c^\infty(\Omega)$ 在 $L^p(\Omega)$ 中稠密，$1\le p<\infty$\\
\texttt{theorem:4-4-2-1} & 近似恒等在 $L^p$ 与 $C_0$ 上的范数恢复\\
\texttt{theorem:4-4-2-2} & 径向可积大函数控制下的 Lebesgue 点恢复\\
\texttt{def:4-4-2-2} & 热核 $p_t(x)=(4\pi t)^{-n/2}e^{-|x|^2/(4t)}$\\
\texttt{def:4-4-3-2} & $C_0$-半群生成元及其定义域\\
\texttt{proposition:4-4-3-1} & $T_tD(A)\subset D(A)$、$AT_tx=T_tAx$ 与轨道微分\\
\texttt{proposition:4-4-3-4} & Laplace--Bochner 预解式、两侧逆、范数界与预解恒等式\\
\bottomrule
\end{longtable}

\section{例题、图示与练习的实施方式}

\subsection{每单元的四拍节奏}

每个单元按下列版式写作，以完整展开关键台阶来形成稳定节奏：
\begin{enumerate}
\item \textbf{可算问题（约 1 页）：}给定具体函数、集合、测度或矩阵，要求计算一个平均、分布函数、变差、卷积或差商。
\item \textbf{定义与最小引理（约 2 页）：}只引入解决该问题所需的术语；每个新符号附一个正例和一个边界例。
\item \textbf{核心证明（约 4--7 页）：}把选择、可测性、交换极限、端点或域假设单列成步骤；每一步标明调用的旧结论。
\item \textbf{回访与变式（约 2 页）：}重新计算开头对象，指出新定理增加了哪项信息，再用一个只改变一项假设的练习收束。
\end{enumerate}

\subsection{例题质量门槛}

每个正文例题同时满足三项：对象完全给定；至少有一个数值、集合、极限或范数可独立核算；所得结论在同单元的定理动机、证明检验或后续回访中出现。纯术语展示改成过渡句；同一个例子再次出现时只计算新的不变量。升级为正文的原练习在章末改成变参数、变假设或构造反例的版本。

建议保留四类图：覆盖选择前后及五倍膨胀；Smith--Volterra 集的拓扑边界与密度窗口；阶梯/Cantor/AC 累计函数及其变差测度；集中、振荡、逃逸三族在空间和尺度上的位置。热核和 Poisson 核使用同一坐标与归一化，图注同时给质量、典型宽度和时间参数。

\subsection{综合练习链}

章末安排四道跨节综合题，每题只组合已在正文证明的工具：
\begin{enumerate}
\item 从 $5r$ 选择、极大弱型到一个指定 $L^1$ 函数的 Lebesgue 点坏集估计；
\item 从混合 Stieltjes 测度到累计函数三分解、普通导数、积分范数与 $N$ 性质；
\item 从弱导数、mollifier 到有界 Lipschitz 域上的 Rellich 子列抽取，并逐项核对 KR 条件；
\item 从 Gaussian 卷积、热核半群到生成元和标量/对角预解式的直接计算。
\end{enumerate}

\section{正文改写顺序与验收标准}

\subsection{实施顺序}

\begin{enumerate}
\item 先更新章首、统一符号、前三章输入卡和标签兼容表；随后锁定四条链的交叉引用。
\item 按 4.1、4.2、4.3、4.4 的依赖次序写作。每完成一个单元即编译，并检查该单元只向前引用定义和定理。
\item 对正式结论建立“陈述--证明--首次调用”三列表；对例题建立“对象--计算--回访”三列表。空缺条目在进入下一大节前补齐。
\item 完成正文后，把原章全部例题、图和练习逐项映射到新位置；正文已经吸收的练习改写成变式。
\item 最后联编第五至七章，核对八个跨章稳定标签及本章内部标签、热核时间归一化以及乘积空间平移、Gaussian 平滑两条新增接口。
\end{enumerate}

\subsection{数学验收清单}

\begin{enumerate}
\item 每个球平均都满足允许半径条件；每个点值式都说明代表；每个上确界或上极限都有可测性来源。
\item centered、uncentered 与含点邻域的约定固定；球、立方体和有界离心率集合的比较写出双向常数。
\item 弱 $L^p$、BV 范数、Stieltjes 端点、分布收敛、Sobolev 域、$C_0$ 强连续和预解集均在首次使用前定义。
\item 4.2 的测度始终定义在 $[a,b]$；左端原子 $c$ 与 $(a,b]$ 增量函数 $F$ 分开记录；任意点态 BV 与规范右连续 BV 的变差区别由孤立尖点例说明。
\item 测度微分到普通差商之间出现有界离心率的测度推论；纯跳跃累计函数的 a.e. 导数通过测度微分识别。
\item 第三章已证的 Riemann 判据、AC 积分表示、$L^1$ 卷积和平移连续作为输入使用；第四章的新增证明具有不同输出。
\item Poincar\'e、延拓、嵌入和 Rellich 的有界性、连通性、Lipschitz 性与指数范围均写入定理陈述；光滑单位分解在使用前构造。
\item KR 的集中族按 $k^{n/p}$ 归一化；平移模负责集中与振荡，尾部条件负责逃逸；临界 Sobolev 例由缩放完成。
\item 近似恒等的质量、统一 $L^1$ 界和尾部集中逐项验证；径向递减大函数只作为逐点恢复的可检验充分条件。
\item 热核满足 $\partial_tp_t=\Delta p_t$；第六章的 $p_{t/2}$ 对应 $\tfrac12\Delta$；Poisson 半群律的证明只调用此前建立的实变量积分。
\item 半群先由强连续和 Baire 得增长界，再定义生成元；Laplace--Bochner 积分限于 $\operatorname{Re}\lambda>\omega$，并证明两个逆恒等式与预解恒等式。
\item 每个 Fubini、DCT、Bochner 积分及极限交换都在相邻文字中给出非负性、绝对可积性或共同支配。
\end{enumerate}

\subsection{结构验收清单}

\begin{enumerate}
\item 四个大节各有一个可计算主问题；4.1、4.2、4.3 各至少有一项输出进入下一大节，4.4 的输出进入第五至七章；章末可用“观测--稳定--极限--识别”复述全部核心定理。
\item 原第四章十二个核心单元均在覆盖表中有落点；合并只消除重复证明与重复计算。
\item 新增正式内容均进入本章核心定理或第五至七章的既有证明；专题边界表中的内容由相应后章或独立发展线完整承载。
\item 正文例题中的常数、支撑与归一化均可逐项核算；每个反例逐项验证保留和破坏的假设。
\item 正文篇幅保持在 215--225 页附近；若版面超出上限，优先合并重复例算与同义图，核心证明的可测性、端点、域条件和极限交换保留完整篇幅。
\item 全卷编译无未解析引用；八个跨章下游标签及本章内部标签仍绑定到相同数学内容；术语首次出现顺序与前三章前置范围一致。
\end{enumerate}

\section{修订后的主线摘要}

第一大节用覆盖几何把全部小尺度的上确界化为弱型估计，再由稠密性恢复函数与测度的局部密度。第二大节把局部密度读成累计函数的普通导数，以端点原子和 Stieltjes 测度保存导数遗漏的跳跃与奇异连续变化，并由 Banach--Zarecki 闭合绝对连续性。第三大节把平移平均写成卷积，以测试函数定义弱导数，再由有限网格、完备性和共同模量得到 Sobolev 紧性。第四大节把重复卷积组织成时间半群，从小时间差商得到生成元，再以指数加权轨道积分获得有界预解式。

因此，第四章的扩充体现为同一证明机器在四种观测中的逐层展开：有限观测给出可计算量，尺度一致估计提供稳定性；稠密逼近、测度正则性、受控收敛、完备性或紧性负责传递并稳定极限；点值、测度、测试函数或生成元完成识别。

\end{document}
